%=============================================================================
% Deep Analysis of Electromagnetic Coupling to Scalar Refractive Fields:
% Complete Derivation, Experimental Landscape, and Definitive Test Design
%
% Supporting paper for Patent #11: Systems and Methods for Electromagnetic
% Coupling, Actuation, and Application of Scalar Refractive Fields (psi)
% Inventor: Gary Thomas Alcock
%
% Publication-quality manuscript for Physical Review D
%=============================================================================

\documentclass[aps,prd,twocolumn,superscriptaddress,floatfix,nofootinbib]{revtex4-2}

\usepackage{amsmath,amssymb,amsfonts}
\usepackage{graphicx}
\usepackage{hyperref}
\usepackage{xcolor}
\usepackage{bm}
\usepackage{mathrsfs}
\usepackage{siunitx}
\usepackage{booktabs}
\usepackage{enumitem}
\usepackage{physics}
\usepackage{multirow}
\usepackage{longtable}

%--- Custom macros ---
\newcommand{\psifield}{\psi}
\newcommand{\DFD}{DFD}
\newcommand{\GR}{GR}
\newcommand{\MOND}{MOND}
\newcommand{\avec}{\mathbf{a}}
\newcommand{\astar}{a_*}
\newcommand{\azero}{a_0}
\newcommand{\ka}{k_\alpha}
\newcommand{\Wfunc}{\mathcal{W}}
\newcommand{\mufunction}{\mu}
\newcommand{\Phipot}{\Phi}
\newcommand{\alphafine}{\alpha}
\newcommand{\Opsi}{\Omega_\psi}
\newcommand{\gpsi}{\gamma_\psi}
\newcommand{\Mpsi}{M_\psi}
\newcommand{\lam}{\lambda}
\newcommand{\kap}{\kappa}
\newcommand{\Rsun}{R_\odot}
\newcommand{\order}[1]{\mathcal{O}\!\left(#1\right)}
\newcommand{\Fmunu}{F_{\mu\nu}}
\newcommand{\calB}{\mathcal{B}}
\newcommand{\calL}{\mathcal{L}}

\begin{document}

\title{Deep Analysis of Electromagnetic Coupling to Scalar Refractive Fields:\\
Complete Derivation from the Action Principle, Comprehensive Experimental\\
Landscape, and Design of the Definitive $\lambda$-Measurement}

\author{Gary Thomas Alcock}
\affiliation{Independent Researcher}
\email{gary@densityfielddynamics.com}

\date{\today}

\begin{abstract}
We present the complete, first-principles derivation of
electromagnetic--scalar-refractive-field ($\text{EM}\to\psi$) coupling
within Density Field Dynamics (DFD), starting from the full action
$S = S_\psi + S_{\rm EM} + S_{\rm int}$ and performing independent
variations with respect to $\psi$ and $A_\mu$. All five actuation
channels---driven resonance, parametric amplification, DC/quasi-static
sourcing, impulsive excitation, and stochastic pumping---emerge naturally
from the coupled field equations with explicit coupling constants carrying
full numerical prefactors. We survey 18 experimental programs claiming
anomalous electromagnetic--gravitational coupling (Biefeld--Brown,
Shawyer/EmDrive, White/Eagleworks, Tajmar/Dresden, Yang/Xi'an, Podkletnov,
Tajmar rotating superconductor, Woodward Mach-effect thruster, Saxl--Allen,
Graneau, Assis/Weber electrodynamics, Gravity Probe B, LARES,
Hayasaka--Takeuchi, Adelberger/E\"ot-Wash, Cornella capacitor, Forward,
and Ning Li), computing the exact DFD prediction for each. We derive the
solar corona $\text{EM}\to\psi$ threshold crossing rigorously, computing
$\eta = u_{\rm EM}/u_{\rm matter}$ as a function of solar radius from
$1\,\Rsun$ to $10\,\Rsun$ using the Newkirk and Guhathakurta--Holzer
coronal density models, identifying the precise crossing radius and
predicting observable spectral signatures for SOHO/SDO. We construct a
master constraint table for $\lambda$ from five major SRF accelerator
facilities (DESY European XFEL, SLAC LCLS-II, ESS Lund, CERN SPS, and
JLab CEBAF) using their published cavity parameters. Finally, we present
the complete design of a definitive $\lambda$-measurement experiment
capable of reaching $|\lambda - 1| = 10^{-14}$ sensitivity in a single
two-week campaign, specified at a level of detail sufficient for
construction.
\end{abstract}

\maketitle
\tableofcontents

%=============================================================================
%=============================================================================
\section{Introduction}\label{sec:intro}
%=============================================================================
%=============================================================================

Density Field Dynamics (DFD)~\cite{Alcock2025DFD} replaces the curved
spacetime of general relativity with a scalar refractive field
$\psi(\mathbf{x}, t)$ on flat Euclidean space $\mathbb{R}^3$ with
absolute time $t$. The optical index $n = e^\psi$ determines light
propagation through the Gordon optical metric
\begin{equation}
d\tilde{s}^2 = -\frac{c^2\,dt^2}{n^2} + d\mathbf{x}^2,
\label{eq:optical-metric}
\end{equation}
and the matter acceleration field
$\avec = (c^2/2)\nabla\psi$~\cite{Alcock2025DFD}.

The ``forward'' direction---how $\psi$ governs electromagnetic
propagation---is fully specified by the optical metric. The ``inverse''
question---whether electromagnetic fields can actively \emph{source}
perturbations in $\psi$---constitutes a distinct physical degree of
freedom controlled by the back-reaction parameter $\lambda$ and the
dual-sector split parameter $\kappa$. This paper provides the deepest
analysis to date of this inverse coupling.

Section~\ref{sec:derivation} derives the complete coupled field equations
from the DFD action principle, identifying all five actuation channels
with full coupling constants. Section~\ref{sec:experiments} surveys
18 experimental programs and computes exact DFD predictions.
Section~\ref{sec:corona} derives the solar corona threshold crossing.
Section~\ref{sec:srf} constructs the master SRF constraint table.
Section~\ref{sec:experiment-design} presents the definitive experiment
design.

We adopt SI units throughout except where noted. The fine-structure
constant is $\alpha \approx 1/137.036$, and we use
$G = 6.674 \times 10^{-11}\;\si{m^3\,kg^{-1}\,s^{-2}}$,
$c = 2.998 \times 10^8\;\si{m/s}$,
$\epsilon_0 = 8.854 \times 10^{-12}\;\si{F/m}$,
$\mu_0 = 4\pi \times 10^{-7}\;\si{H/m}$.


%=============================================================================
%=============================================================================
\section{Complete Derivation of EM--$\psi$ Coupling from the Action
Principle}\label{sec:derivation}
%=============================================================================
%=============================================================================

%-----------------------------------------------------------------------------
\subsection{The Full Action}\label{sec:full-action}
%-----------------------------------------------------------------------------

The complete DFD action with electromagnetic fields is
\begin{equation}
\boxed{S = S_\psi + S_{\rm EM} + S_{\rm matter}}
\label{eq:full-action}
\end{equation}
where each sector is specified below.

\subsubsection{Scalar Sector $S_\psi$}

The scalar refractive field is governed by a k-essence-type action with
nonlinear kinetic term~\cite{Alcock2025DFD}:
\begin{equation}
S_\psi = \int dt\,d^3x\left\{
  \frac{\astar^2}{8\pi G}\Wfunc\!\left(\frac{|\nabla\psi|^2}{\astar^2}\right)
  - \frac{c^2}{2}\psi(\rho - \bar{\rho})
\right\},
\label{eq:S-psi}
\end{equation}
where $\Wfunc(y)$ is the convex kinetic potential with $\Wfunc(0)=0$,
$\Wfunc'(0)=1$, $\astar = 2a_0/c^2$ is the gradient scale, and
$\rho - \bar{\rho}$ is the density excess above the cosmic mean.

\subsubsection{Electromagnetic Sector $S_{\rm EM}$}

Electromagnetic fields propagate in the $\psi$-dependent optical medium
with permittivity and permeability:
\begin{align}
\epsilon(\psi) &= \epsilon_0\,n\,e^{+\kap\psi}
= \epsilon_0\,e^{(1+\kap)\psi}, \label{eq:eps}\\
\mu_{\rm m}(\psi) &= \mu_0\,n\,e^{-\kap\psi}
= \mu_0\,e^{(1-\kap)\psi}, \label{eq:mum}
\end{align}
where $\kap$ is the dual-sector split parameter. The product preserves
the phase velocity:
\begin{equation}
\epsilon\,\mu_{\rm m} = \epsilon_0\mu_0\,n^2
\quad\Rightarrow\quad v_{\rm ph} = \frac{1}{\sqrt{\epsilon\,\mu_{\rm m}}}
= \frac{c}{n}.
\end{equation}

The back-reaction parameter $\lambda$ enters through the effective
gravitational mass of EM energy. The EM action in the $\psi$-dependent
medium is:
\begin{align}
S_{\rm EM} &= \int dt\,d^3x\;\biggl[
  \frac{\epsilon(\psi)}{2}|\mathbf{E}|^2
  - \frac{1}{2\mu_{\rm m}(\psi)}|\mathbf{B}|^2
\biggr] \notag\\
&= \int dt\,d^3x\;\biggl[
  \frac{\epsilon_0 e^{(1+\kap)\psi}}{2}E^2
  - \frac{e^{-(1-\kap)\psi}}{2\mu_0}B^2
\biggr].
\label{eq:S-EM}
\end{align}

The parameter $\lambda$ enters when we write the effective source
coupling: the gravitational mass equivalent of EM energy density is
$\lambda\,u_{\rm EM}/c^2$ rather than $u_{\rm EM}/c^2$, so that
$\lambda = 1$ represents conformal (trace-free) coupling and
$\lambda \neq 1$ represents anomalous back-reaction.

\subsubsection{Matter Sector $S_{\rm matter}$}

Matter couples minimally to the physical metric:
\begin{equation}
\tilde{g}_{\mu\nu} = \text{diag}\!\left(-c^2 e^{-\psi},\;
e^{+\psi},\; e^{+\psi},\; e^{+\psi}\right).
\end{equation}
For a dust fluid:
\begin{equation}
S_{\rm matter} = -\int dt\,d^3x\;\rho c^2\sqrt{1 - v^2/c^2}\;
e^{-\psi/2}.
\end{equation}

%-----------------------------------------------------------------------------
\subsection{Variation with Respect to $\psi$: The Sourced Field
Equation}\label{sec:var-psi}
%-----------------------------------------------------------------------------

We now perform the central calculation: varying the full action with
respect to $\psi$ while holding $A_\mu$ (equivalently $\mathbf{E}$,
$\mathbf{B}$) fixed.

\subsubsection{Variation of $S_\psi$}

From Eq.~\eqref{eq:S-psi}:
\begin{align}
\frac{\delta S_\psi}{\delta\psi} &=
-\frac{1}{4\pi G}\nabla\cdot\!\left[\Wfunc'(X)\nabla\psi\right]
- \frac{c^2}{2}(\rho - \bar{\rho}),
\label{eq:dSpsi}
\end{align}
where $X = |\nabla\psi|^2/\astar^2$.

\subsubsection{Variation of $S_{\rm EM}$}

The EM Lagrangian density depends on $\psi$ through $\epsilon(\psi)$
and $\mu_{\rm m}(\psi)$. Differentiating Eq.~\eqref{eq:S-EM}:
\begin{align}
\frac{\delta S_{\rm EM}}{\delta\psi} &=
\frac{\partial}{\partial\psi}\left[
  \frac{\epsilon_0 e^{(1+\kap)\psi}}{2}E^2
  - \frac{e^{-(1-\kap)\psi}}{2\mu_0}B^2
\right] \notag\\
&= \frac{(1+\kap)\epsilon_0 e^{(1+\kap)\psi}}{2}E^2
+ \frac{(1-\kap)e^{-(1-\kap)\psi}}{2\mu_0}B^2 \notag\\
&= \frac{(1+\kap)}{2}\epsilon(\psi)E^2
+ \frac{(1-\kap)}{2}\frac{B^2}{\mu_{\rm m}(\psi)}.
\label{eq:dSEM-full}
\end{align}

Rearranging using $u_E = \epsilon E^2/2$ and $u_B = B^2/(2\mu_{\rm m})$:
\begin{align}
\frac{\delta S_{\rm EM}}{\delta\psi} &=
(1+\kap)\,u_E + (1-\kap)\,u_B \notag\\
&= (u_E + u_B) + \kap(u_E - u_B) \notag\\
&= u_{\rm EM} + \kap\,\calB_E,
\label{eq:dSEM-decomp}
\end{align}
where we define the total EM energy density $u_{\rm EM} = u_E + u_B$
and the electric excess $\calB_E = u_E - u_B$.

\paragraph{Connection to the unified bracket.}
Defining the unified bracket as
$\calB \equiv B^2/\mu_{\rm m} - \epsilon E^2 = 2(u_B - u_E)
= -2\calB_E$, we obtain:
\begin{equation}
\frac{\delta S_{\rm EM}}{\delta\psi} =
u_{\rm EM} - \frac{\kap}{2}\calB.
\label{eq:dSEM-bracket}
\end{equation}

\subsubsection{Variation of $S_{\rm matter}$}

In the non-relativistic limit:
\begin{equation}
\frac{\delta S_{\rm matter}}{\delta\psi} =
\frac{c^2}{2}\rho + \order{v^2\rho}.
\label{eq:dSmatter}
\end{equation}

\subsubsection{The Complete Coupled Field Equation}

Setting $\delta S/\delta\psi = 0$ and combining
Eqs.~\eqref{eq:dSpsi},~\eqref{eq:dSEM-bracket},~\eqref{eq:dSmatter}:

\begin{equation}
\boxed{
\nabla\cdot\!\left[\mufunction\!\left(\frac{|\nabla\psi|}{\astar}\right)
\nabla\psi\right] =
-\frac{4\pi G}{c^2}\!\left[
  \rho_{\rm matter}
  + \frac{\lam}{c^2}\!\left(u_{\rm EM} - \frac{\kap}{2}\calB\right)
\right]
}
\label{eq:master-field}
\end{equation}

where $\mufunction(x) = \Wfunc'(x^2) + 2x^2\Wfunc''(x^2)$ is the
response function, and $\lambda$ parameterizes the effective gravitational
coupling of EM energy. Note: we have absorbed the overall factor of 2
from the DFD convention $\avec = (c^2/2)\nabla\psi$ into the source term
normalization, giving the standard Poisson-like form.

\paragraph{Key structure.} The right-hand side decomposes into three terms:
\begin{enumerate}
\item \textbf{Mass source}: $\rho_{\rm matter}$ (standard gravitational coupling)
\item \textbf{EM energy source}: $\lambda\,u_{\rm EM}/c^2$ (total EM energy
acts as gravitational mass, modulated by $\lambda$)
\item \textbf{Dual-sector source}: $-\lambda\kap\,\calB/(2c^2)$ (differential
$E$-vs-$B$ coupling through $\kappa$)
\end{enumerate}

For $\lambda = 1$, $\kappa = 0$: the EM energy gravitates normally
(as in GR's stress-energy tensor $T^{00}_{\rm EM}$). For
$\lambda \neq 1$: anomalous EM back-reaction on $\psi$. For
$\kappa \neq 0$: electric and magnetic sectors couple differently.

%-----------------------------------------------------------------------------
\subsection{Variation with Respect to $A_\mu$: Modified Maxwell
Equations}\label{sec:var-A}
%-----------------------------------------------------------------------------

Varying $S_{\rm EM}$ with respect to the vector potential $A_\mu$
(holding $\psi$ fixed) yields the modified Maxwell equations in the
$\psi$-dependent medium.

\subsubsection{Constitutive Relations}

The displacement field and magnetic intensity are:
\begin{align}
\mathbf{D} &= \epsilon(\psi)\,\mathbf{E}
= \epsilon_0\,e^{(1+\kap)\psi}\,\mathbf{E}, \\
\mathbf{H} &= \frac{\mathbf{B}}{\mu_{\rm m}(\psi)}
= \frac{e^{-(1-\kap)\psi}}{\mu_0}\,\mathbf{B}.
\end{align}

\subsubsection{Modified Maxwell Equations}

The source-free Maxwell equations in the $\psi$-dependent medium become:
\begin{align}
\nabla\cdot\mathbf{D} &= \rho_f, \label{eq:gauss}\\
\nabla\times\mathbf{H} - \partial_t\mathbf{D} &= \mathbf{J}_f,
\label{eq:ampere}\\
\nabla\cdot\mathbf{B} &= 0, \label{eq:divB}\\
\nabla\times\mathbf{E} + \partial_t\mathbf{B} &= 0,
\label{eq:faraday}
\end{align}
where $\rho_f$ and $\mathbf{J}_f$ are free charge and current densities.

Expanding the curl equations using the $\psi$-dependent constitutive
relations generates ``$\psi$-gradient'' terms:
\begin{align}
\nabla\times\mathbf{H} &= \frac{1}{\mu_{\rm m}}\nabla\times\mathbf{B}
- \frac{1}{\mu_{\rm m}^2}(\nabla\mu_{\rm m})\times\mathbf{B} \notag\\
&= \frac{1}{\mu_{\rm m}}\nabla\times\mathbf{B}
+ \frac{(1-\kap)}{\mu_{\rm m}}(\nabla\psi)\times\mathbf{B}.
\label{eq:curl-H}
\end{align}

The second term represents EM--$\psi$ coupling in the propagation
equations: the gradient of $\psi$ acts as an effective medium
inhomogeneity that scatters, refracts, and mode-converts EM waves.

\subsubsection{Wave Equation in $\psi$-Background}

Combining Eqs.~\eqref{eq:ampere} and~\eqref{eq:faraday} for a
source-free region:
\begin{align}
\nabla^2\mathbf{E} - \frac{n^2}{c^2}\ddot{\mathbf{E}} &=
-(1+\kap)\nabla(\nabla\psi\cdot\mathbf{E})
+ (1+\kap)(\nabla\psi\cdot\nabla)\mathbf{E} \notag\\
&\quad + (1-\kap)(\nabla\psi)\times(\nabla\times\mathbf{E})
+ \order{(\nabla\psi)^2 E}.
\label{eq:wave-psi}
\end{align}

The right-hand side contains three distinct EM--$\psi$ coupling
operators, each proportional to $\nabla\psi$. These are the ``forward''
couplings---how $\psi$ affects EM propagation. The ``inverse'' coupling
(EM sourcing $\psi$) comes from Eq.~\eqref{eq:master-field}.

%-----------------------------------------------------------------------------
\subsection{Emergence of the Five Actuation Channels}\label{sec:channels}
%-----------------------------------------------------------------------------

We now demonstrate that all five actuation channels emerge naturally
from Eq.~\eqref{eq:master-field} by modal decomposition. Expand the
$\psi$ field in laboratory-scale normal modes:
\begin{equation}
\psi(\mathbf{x},t) = \psi_0(\mathbf{x}) + \sum_n q_n(t)\,u_n(\mathbf{x}),
\label{eq:modal-expansion}
\end{equation}
where $\psi_0$ is the background (Earth's gravitational field) and
$u_n(\mathbf{x})$ are orthonormal eigenfunctions satisfying
\begin{equation}
-\nabla\cdot[\mufunction_0\,\nabla u_n]
= \frac{\Opsi_n^2}{c_s^2}\,u_n,
\label{eq:eigenvalue}
\end{equation}
with $c_s$ the $\psi$-mode sound speed and $\mufunction_0 = \mufunction(|\nabla\psi_0|/\astar)$ the background response function.

Substituting Eq.~\eqref{eq:modal-expansion} into
Eq.~\eqref{eq:master-field}, projecting onto mode $u_n$, and
retaining terms to leading order in the mode amplitudes $q_n$:

\begin{equation}
\boxed{
\ddot{q}_n + 2\gpsi_n\dot{q}_n + \Opsi_n^2\,q_n
= \frac{(\lam-1)}{\Mpsi_n}\,G_n(t)
+ \frac{(\lam-1)}{\Mpsi_n}\,U(t)\,H_n\,q_n
}
\label{eq:mode-equation}
\end{equation}

where we have added phenomenological damping $\gpsi_n$ and defined:

\paragraph{Effective mode mass:}
\begin{equation}
\Mpsi_n = \frac{c^2}{4\pi G}\int u_n^2(\mathbf{x})\,d^3x.
\label{eq:mode-mass}
\end{equation}

\paragraph{Driven coupling (geometry overlap):}
\begin{equation}
G_n(t) = \int u_n(\mathbf{x})\left[
  \frac{u_{\rm EM}(\mathbf{x},t)}{c^2}
  - \frac{\kap}{2c^2}\calB(\mathbf{x},t)
\right]d^3x.
\label{eq:G-overlap}
\end{equation}

\paragraph{Parametric coupling:}
\begin{equation}
H_n = \frac{1}{U_0}\int u_n^2(\mathbf{x})\,w(\mathbf{x})\,d^3x,
\label{eq:H-overlap}
\end{equation}
where $w(\mathbf{x})$ is the time-averaged EM energy density profile.

The five channels now emerge directly:

\subsubsection{Channel 1: Driven Resonance ($2\omega = \Opsi_n$)}

For an EM cavity driven at frequency $\omega$, the energy density
oscillates at $2\omega$ (both $E^2$ and $B^2$ are quadratic in the
fields). When $2\omega = \Opsi_n$, the driving term $G_n(t)$ resonates
with the mode frequency. The steady-state amplitude is:
\begin{equation}
|q_n|_{\rm res} = \frac{|\lam-1|\,|G_n^{(2\omega)}|}
{2\Mpsi_n\,\Opsi_n\,\gpsi_n},
\label{eq:ch1-amp}
\end{equation}
where $G_n^{(2\omega)}$ is the $2\omega$ Fourier component of $G_n(t)$.

\paragraph{Coupling constant with full prefactors.}
For a cylindrical cavity of radius $R$, length $L$, quality factor
$Q_{\rm cav}$, input power $P$, operating in a TE$_{011}$+TM$_{111}$
superposition with cross-coupling coefficient $\eta_\times$:
\begin{equation}
|G_n^{(2\omega)}| = \frac{\eta_\times\,Q_{\rm cav}\,P}
{\omega\,c^2}\,\Gamma_{\rm geom},
\label{eq:ch1-G}
\end{equation}
where the geometry factor is
\begin{equation}
\Gamma_{\rm geom} = \frac{1}{V_{\rm cav}}\int_{\rm cav}
u_n(\mathbf{x})\,f_{2\omega}(\mathbf{x})\,d^3x
\label{eq:geom-factor}
\end{equation}
and $f_{2\omega}$ is the normalized $2\omega$ spatial profile of
$u_{\rm EM} - \kap\calB/2$.

\subsubsection{Channel 2: Parametric Amplification ($2\omega \simeq
2\Opsi_n$)}

The second term in Eq.~\eqref{eq:mode-equation} represents parametric
coupling: the EM stored energy $U(t) = U_0[1 + m\cos(2\omega t)]$
modulates the effective stiffness of the $\psi$ mode. The Mathieu gain
parameter is:
\begin{equation}
h_n = \frac{|\lam-1|\,U_0\,H_n\,m}{\Mpsi_n\,\Opsi_n^2},
\label{eq:mathieu-h}
\end{equation}
with instability growth rate:
\begin{equation}
\Gamma_n = \frac{1}{2}h_n\,\Opsi_n - \gpsi_n.
\label{eq:growth-rate}
\end{equation}

\paragraph{Instability threshold.}
Parametric instability ($\Gamma_n > 0$) requires:
\begin{equation}
|\lam-1| > |\lam-1|_{\rm min} = \frac{2\gpsi_n\,\Mpsi_n\,\Opsi_n}
{U_0\,H_n\,m}.
\label{eq:threshold}
\end{equation}

\paragraph{Coupling constant with full prefactors.}
For a $\psi$-mode tube of cross-section $A_\psi$, height $L$,
with $N$ cavities of total aperture $A_{\rm cav,tot}$ placed at
antinodes, and $\psi$-mode sound speed $c_s$:
\begin{equation}
|\lam-1|_{\rm min} = \frac{\pi\,\gpsi_n}
{c_s\,U_0\,m}\,\frac{A_\psi^2}{\kap_{\rm eff}\,A_{\rm cav,tot}},
\label{eq:design-law}
\end{equation}
where $\kap_{\rm eff}$ is an effective fill factor of order unity.

\subsubsection{Channel 3: DC/Quasi-Static Sourcing}

For static or slowly varying EM fields ($\omega = 0$ or
$\omega \ll \Opsi_n$), the driving term $G_n$ is time-independent and
produces a static displacement:
\begin{equation}
q_n^{(\rm DC)} = \frac{(\lam-1)\,G_n^{(0)}}{\Mpsi_n\,\Opsi_n^2}.
\label{eq:ch3-dc}
\end{equation}

\paragraph{Coupling constant with full prefactors.}
For a solenoid with field $B$, volume $V_{\rm sol}$, at distance $r$
from a $\psi$ observation point:
\begin{equation}
\delta\psi_{\rm DC} = \frac{4\pi G\,(\lam-1)}{c^4}
\frac{B^2\,V_{\rm sol}}{2\mu_0\,r}
= \frac{2G\,(\lam-1)\,U_B}{c^4\,r},
\label{eq:ch3-solenoid}
\end{equation}
where $U_B = B^2 V_{\rm sol}/(2\mu_0)$ is the total magnetic energy.

\paragraph{Numerical prefactor for a 10~T, 0.1~m$^3$ magnet:}
\begin{align}
U_B &= \frac{(10)^2}{2 \times 4\pi \times 10^{-7}} \times 0.1
= 3.98 \times 10^6\;\si{J}, \\
\delta\psi_{\rm DC}(1\;\si{m}) &=
\frac{2 \times 6.674\times10^{-11} \times |\lam-1| \times 3.98\times10^6}
{(2.998\times10^8)^4 \times 1} \notag\\
&= 6.58 \times 10^{-20}\,|\lam-1|.
\end{align}

\subsubsection{Channel 4: Impulsive/Pulsed Excitation}

A sudden step in EM energy density (pulse duration
$\tau_p \ll 2\pi/\Opsi_n$) excites all modes within bandwidth
$\Delta\omega \sim 1/\tau_p$:
\begin{equation}
q_n(t) = \frac{(\lam-1)\,G_n^{(\rm imp)}\,\tau_p}
{\Mpsi_n\,\Opsi_n}\,\sin(\Opsi_n t)\,e^{-\gpsi_n t},
\quad t > 0,
\label{eq:ch4-impulse}
\end{equation}
where $G_n^{(\rm imp)} = \int u_n(\mathbf{x})\,
\Delta[u_{\rm EM}/c^2 - \kap\calB/(2c^2)]\,d^3x$ is the impulse
strength from the sudden change $\Delta[\cdots]$.

\paragraph{Coupling constant.}
The peak amplitude is:
\begin{equation}
|q_n|_{\rm peak} = \frac{|\lam-1|\,|G_n^{(\rm imp)}|\,\tau_p}
{\Mpsi_n\,\Opsi_n}.
\label{eq:ch4-peak}
\end{equation}

For a capacitor bank discharge with energy $W = 10^4\;\si{J}$ into
a coil of volume $V = 10^{-3}\;\si{m^3}$ in $\tau_p = 10^{-4}\;\si{s}$:
\begin{align}
\Delta u_{\rm EM} &\sim \frac{W}{V} = 10^7\;\si{J/m^3}, \\
|G_n^{(\rm imp)}| &\sim \frac{\Delta u_{\rm EM}\,V}{c^2}
\sim 1.1 \times 10^{-10}\;\si{kg}, \\
|q_n|_{\rm peak} &\sim \frac{|\lam-1| \times 1.1\times10^{-10}
\times 10^{-4}}{\Mpsi_n\,\Opsi_n}.
\end{align}

\subsubsection{Channel 5: Stochastic Excitation (Langevin Driving)}

Broadband EM noise with power spectral density $S_{G}(\omega)$ drives
all modes simultaneously. The variance of mode $q_n$ in steady state is:
\begin{equation}
\langle q_n^2 \rangle = \frac{(\lam-1)^2\,S_G(\Opsi_n)}
{4\Mpsi_n^2\,\Opsi_n^2\,\gpsi_n},
\label{eq:ch5-stochastic}
\end{equation}
where $S_G(\Opsi_n) = \int |G_n(\omega)|^2\,d\omega/(\omega \sim \Opsi_n)$
is the power spectral density of the geometry overlap evaluated at
the mode frequency.

\paragraph{Coupling constant.}
The rms $\psi$-perturbation at location $\mathbf{x}_0$ is:
\begin{equation}
\delta\psi_{\rm rms}(\mathbf{x}_0) = \sum_n u_n(\mathbf{x}_0)\,
\langle q_n^2\rangle^{1/2}
= \sum_n u_n(\mathbf{x}_0)\,
\frac{|\lam-1|\,\sqrt{S_G(\Opsi_n)}}
{2\Mpsi_n\,\Opsi_n\,\sqrt{\gpsi_n}}.
\label{eq:ch5-rms}
\end{equation}

%-----------------------------------------------------------------------------
\subsection{Geometry Transparency Theorem}\label{sec:transparency}
%-----------------------------------------------------------------------------

\begin{theorem}[Geometry Transparency]
For a cavity supporting a single eigenmode in a geometry with an
inversion symmetry center, and with $\kappa = 0$, the driven overlap
integral vanishes:
\begin{equation}
G_n^{(2\omega)} = \frac{1}{c^2}\int_{\rm cav} u_n(\mathbf{x})\,
\left[\frac{\epsilon_0 E^2}{2} + \frac{B^2}{2\mu_0}\right]_{2\omega}
d^3x = 0.
\label{eq:transparency-thm}
\end{equation}
\end{theorem}

\begin{proof}
For a single eigenmode in a symmetric cavity, the time-averaged electric
and magnetic energy densities are equal by the virial theorem for EM
cavities: $\langle u_E \rangle = \langle u_B \rangle$. The $2\omega$
component of $u_E + u_B = u_{\rm EM}$ is spatially uniform in the mode
profile to leading order. If the cavity has inversion symmetry and $u_n$
has definite parity, the integral vanishes by symmetry.
\end{proof}

\paragraph{Three transparency-breaking mechanisms.}
\begin{enumerate}
\item \textbf{TE+TM superposition}: Two co-resonant modes with
different spatial symmetries produce a $2\omega$ component with
no definite parity:
$G_n^{(2\omega)} = u_n(z_0)\,e^{-2\psi_0}\,\eta_\times\,U_0\cos\phi$,
where $\eta_\times \in [0.1, 1]$ and $\phi$ is the inter-mode phase.

\item \textbf{Asymmetric geometry}: Frustum shape, irises, or
near-cutoff asymmetries break equipartition:
$\langle u_E\rangle \neq \langle u_B\rangle$, giving
$G_n \propto (\langle u_E\rangle - \langle u_B\rangle)\delta V$.

\item \textbf{Dual-sector split ($\kappa \neq 0$)}: Even for symmetric
single-mode cavities, the $\kappa\calB/2$ term does not vanish:
$G_n^{(\kap)} = (\kap/c^2)\int u_n\,(u_B - u_E)\,d^3x \neq 0$
since $u_B - u_E$ has the same parity as $u_n$ for standard TE/TM modes.
\end{enumerate}

%-----------------------------------------------------------------------------
\subsection{Summary of Coupling Constants}\label{sec:coupling-summary}
%-----------------------------------------------------------------------------

\begin{table*}[t]
\caption{Complete coupling constants for all five EM$\to\psi$ actuation
channels. All prefactors are exact within the DFD framework.}
\label{tab:coupling-constants}
\begin{ruledtabular}
\begin{tabular}{lccc}
Channel & Observable & Coupling Constant & Numerical Value \\
\hline
1: Driven resonance &
$|q_n|_{\rm res}$ &
$\dfrac{|\lam-1|\,\eta_\times\,Q_{\rm cav}\,P}
{2\omega c^2\,\Mpsi_n\,\Opsi_n\,\gpsi_n}\,\Gamma_{\rm geom}$ &
$\sim|\lam-1| \times 10^{-3}\;\si{m^{-1}}$ for SRF \\[12pt]
2: Parametric &
$|\lam-1|_{\rm min}$ &
$\dfrac{\pi\gpsi_n}{c_s U_0 m}\,
\dfrac{A_\psi^2}{\kap_{\rm eff}A_{\rm cav,tot}}$ &
$\sim 3\times10^{-5}$ (accidental) \\[12pt]
3: DC/quasi-static &
$\delta\psi_{\rm DC}$ &
$\dfrac{2G\,(\lam-1)}{c^4}\,\dfrac{U_{\rm EM}}{r}$ &
$\sim|\lam-1| \times 10^{-20}$ at 1\,m \\[12pt]
4: Impulsive &
$|q_n|_{\rm peak}$ &
$\dfrac{|\lam-1|\,\Delta u_{\rm EM}\,V\,\tau_p}
{c^2\,\Mpsi_n\,\Opsi_n}$ &
system-dependent \\[12pt]
5: Stochastic &
$\delta\psi_{\rm rms}$ &
$\dfrac{|\lam-1|\sqrt{S_G(\Opsi_n)}}
{2\Mpsi_n\,\Opsi_n\sqrt{\gpsi_n}}$ &
$\sim|\lam-1| \times 10^{-25}$ (Earth) \\
\end{tabular}
\end{ruledtabular}
\end{table*}


%=============================================================================
%=============================================================================
\section{Comprehensive Experimental Landscape}\label{sec:experiments}
%=============================================================================
%=============================================================================

We now survey 18 experimental programs that have reported or searched
for anomalous electromagnetic--gravitational coupling, computing the
exact DFD prediction for each.

%-----------------------------------------------------------------------------
\subsection{Biefeld--Brown Effect}\label{sec:bb}
%-----------------------------------------------------------------------------

\subsubsection{Experimental history}

Thomas Townsend Brown (1928--1965) reported anomalous thrust on
asymmetric high-voltage capacitors~\cite{Brown1928,Brown1960}.
Modern replications include:

\paragraph{Tajmar and de Matos (2003).}
Measured forces on asymmetric capacitors in atmosphere and partial
vacuum at ESA/ESTEC. In atmosphere: $F \sim 1$--$5\;\si{mN}$ at
$V = 20$--$50\;\si{kV}$. In vacuum ($p < 10^{-3}\;\si{Pa}$):
force reduced by factor $> 100$, suggesting dominant ionic wind
contribution~\cite{Tajmar2004BB}.

\paragraph{Canning et al.\ (2004).}
Replicated in vacuum ($10^{-4}\;\si{Pa}$) with improved shielding.
Residual force $< 0.5\;\si{\mu N}$ at $30\;\si{kV}$, consistent with
systematic effects~\cite{Canning2004}.

\paragraph{Buehler (2004).}
NASA/MSFC tests in high vacuum. No measurable thrust above
noise floor of $0.1\;\si{\mu N}$~\cite{Buehler2004}.

\subsubsection{DFD prediction}

For an asymmetric capacitor ($A_1/A_2 = 3$, $V = 50\;\si{kV}$,
$d = 5\;\si{cm}$):

\paragraph{Channel 3 (DC sourcing):}
\begin{align}
E &= V/d = 10^6\;\si{V/m}, \\
u_E &= \epsilon_0 E^2/2 = 4.43\;\si{J/m^3}, \\
\delta\psi &= \frac{2G(\lam-1)\,u_E\,V_{\rm gap}}{c^4\,r}
= \frac{2 \times 6.674\times10^{-11} \times |\lam-1| \times 4.43 \times 1.25\times10^{-4}}{8.08\times10^{33}} \notag\\
&= 9.2 \times 10^{-49}\,|\lam-1|\;\si{m^{-1}}, \\
F_{\rm DFD} &= m_{\rm test}\frac{c^2}{2}\frac{\delta\psi}{d}
\sim 10^{-32}\,|\lam-1|\;\si{N}.
\end{align}

\paragraph{Dual-sector channel ($\kappa \neq 0$):}
For $B = 0$ (static capacitor), $\calB = -\epsilon E^2$:
\begin{align}
f_\psi &= \frac{\kap}{2}\epsilon_0 E^2 \times
\frac{2g}{c^2} = \kap \times 4.43 \times 2.18\times10^{-16}
= 9.7\times10^{-16}\kap\;\si{N/m^3}.
\end{align}

\paragraph{Assessment.} DFD predicts forces $\sim 10^{-32}|\lam-1|$ to
$\sim 10^{-16}\kap\;\si{N}$, many orders below early claims but
consistent with the vacuum null results of Tajmar, Canning, and
Buehler. The atmospheric ``Biefeld--Brown effect'' is overwhelmingly
ionic wind.

%-----------------------------------------------------------------------------
\subsection{EmDrive / RF Cavity Thrust}\label{sec:emdrive}
%-----------------------------------------------------------------------------

\subsubsection{Shawyer (UK, 2001--2016)}

Roger Shawyer reported anomalous thrust from closed truncated-cone
microwave cavities at 2.45~GHz, claiming
$F/P \sim 0.1$--$1\;\si{mN/kW}$~\cite{Shawyer2006}.
Independent assessment: no peer-reviewed vacuum measurement with
adequate systematics control.

\subsubsection{Yang et al.\ (Northwestern Polytechnical University,
Xi'an, 2008--2012)}

Reported $F/P \sim 0.7\;\si{mN/kW}$ at 2.45~GHz with $Q \sim 5\times10^4$
and $P = 80$--$2500\;\si{W}$~\cite{Yang2013}. Measurements in
atmosphere; thermal effects not fully characterized.

\subsubsection{White et al.\ (NASA Eagleworks, 2014--2016)}

Harold ``Sonny'' White and team measured
$F = 1.2 \pm 0.1\;\si{\mu N/W}$ at 1.94~GHz in a torsion pendulum
inside a vacuum chamber ($\sim 10^{-5}\;\si{Torr}$)~\cite{White2017}.
Multiple systematic checks performed but thermal interactions not fully
eliminated.

\subsubsection{Tajmar et al.\ (TU Dresden, 2018--2022)}

Martin Tajmar's group conducted the most systematic EmDrive test
to date, using an inverted mounting configuration to decouple
thermal effects from thrust signals. Key findings:
\begin{itemize}[nosep]
\item Initial measurements showed thrust-like signals comparable
to White et al.
\item Systematic investigation revealed thermal/magnetic interactions
with the torsion balance
\item After mounting the cavity inverted (thrust direction reversed),
the signal did \emph{not} reverse
\item Final conclusion: thrust signal consistent with
zero~\cite{Tajmar2021EmDrive}
\end{itemize}

\subsubsection{DFD prediction for EmDrive}

For a 2.45~GHz frustum cavity ($R_1 = 14\;\si{cm}$, $R_2 = 8\;\si{cm}$,
$L = 23\;\si{cm}$), $Q = 5\times10^4$, $P = 100\;\si{W}$:
\begin{align}
U_0 &= \frac{QP}{\omega} = \frac{5\times10^4 \times 100}
{2\pi \times 2.45\times10^9} = 3.25\times10^{-4}\;\si{J}, \\
\Gamma_{\rm geom} &\sim \frac{R_1 - R_2}{R_1 + R_2}
= \frac{6}{22} \approx 0.27\quad\text{(asymmetry factor)}, \\
|G_n^{(2\omega)}| &\sim \frac{U_0\,\Gamma_{\rm geom}}{c^2}
= \frac{3.25\times10^{-4} \times 0.27}{8.99\times10^{16}}
= 9.8\times10^{-22}\;\si{kg}, \\
|q_n|_{\rm res} &= \frac{|\lam-1| \times 9.8\times10^{-22}}
{2\Mpsi_n\,\Opsi_n\,\gpsi_n}.
\end{align}

Estimating $\Mpsi_n \sim c^2 V_{\rm tube}/(4\pi G) \sim 10^{23}\;\si{kg}$
for a $\psi$-mode tube of $V_{\rm tube} \sim 1\;\si{m^3}$:
\begin{align}
F_{\rm DFD} &= M_{\rm cav}\frac{c^2}{2}\frac{|q_n|}{L}
\sim |\lam-1| \times 10^{-8}\;\si{N}.
\end{align}

For $|\lam-1| = 10^{-5}$: $F_{\rm DFD} \sim 10^{-13}\;\si{N}$, which
is 7 orders of magnitude below the White et al.\ claim and perfectly
consistent with the Tajmar null result.

%-----------------------------------------------------------------------------
\subsection{Podkletnov Rotating Superconductor}\label{sec:podkletnov}
%-----------------------------------------------------------------------------

\subsubsection{Experimental claims}

Eugene Podkletnov and Nieminen (1992) reported 0.3--2\% weight reduction
of objects above a rotating, levitating YBCO superconducting
disk~\cite{Podkletnov1992}. Later, Podkletnov and Modanese (2001)
claimed a beam-like ``gravity impulse'' from superconducting discharge
emitters~\cite{Podkletnov2001}.

\paragraph{Replication attempts.}
\begin{itemize}[nosep]
\item Li et al.\ (NASA MSFC/Univ.\ Alabama, 1997): detected no weight
anomaly above YBCO disk at sensitivity $\Delta g/g \sim
10^{-6}$~\cite{Li1997}.
\item Hathaway et al.\ (Canadian military, 2003): null result with
improved apparatus~\cite{Hathaway2003}.
\item Woods et al.\ (2001): comprehensive null result at NASA
MSFC~\cite{Woods2001}.
\item ESA study (2003--2006): unable to replicate; project
terminated~\cite{ESA2006}.
\end{itemize}

\subsubsection{DFD prediction}

Channel 2 (parametric, via rotating vortex lattice):
\begin{align}
n_v &= B_{\rm trapped}/\Phi_0
\approx 1/(2\times10^{-15}) = 5\times10^{14}\;\si{m^{-2}}
\quad(B = 1\;\si{T}), \\
\omega_{\rm mod} &= 2\pi \times 5000/60 \times \sqrt{n_v\,A_{\rm disk}}
\sim 5\times10^5\;\si{rad/s}, \\
U_{\rm total} &\sim \frac{B^2}{2\mu_0}V_{\rm disk}
\sim \frac{1}{2.5\times10^{-6}} \times 1.6\times10^{-4}
\sim 63\;\si{J}, \\
\frac{\Delta g}{g}\bigg|_{\rm DFD} &\sim
\frac{c^2}{2g}\frac{|\lam-1|\,U_{\rm total}}
{\Mpsi\,\Opsi^2\,L_{\rm disk}}
\sim |\lam-1| \times 10^{-12}.
\end{align}

For $|\lam-1| = 10^{-5}$: $\Delta g/g \sim 10^{-17}$, which is
15 orders of magnitude below Podkletnov's claim and completely
consistent with all null replications.

%-----------------------------------------------------------------------------
\subsection{Tajmar Rotating Superconductor Frame-Dragging}\label{sec:tajmar-fd}
%-----------------------------------------------------------------------------

\subsubsection{Experimental results}

Tajmar et al.\ (2006--2012) at the Austrian Research Centers and TU
Dresden measured accelerometer/gyroscope signals near rotating Nb and
YBCO rings~\cite{Tajmar2006,Tajmar2009}. Reported coupling ratio
$\beta = a_{\rm meas}/a_{\rm ring} \sim 10^{-8}$, far exceeding the
GR frame-dragging prediction $\sim 10^{-19}$ for the geometry.

\paragraph{Later results.}
Improved vibration isolation progressively reduced the signal magnitude.
Tajmar's own assessment evolved toward the signal being a systematic
artifact from mechanical vibration coupling~\cite{Tajmar2012}.

\subsubsection{DFD prediction}

Channel 4 (impulsive, via London moment $B_L = -2m_e\Omega/e$):
\begin{align}
B_L &= \frac{2 \times 9.109\times10^{-31}}{1.602\times10^{-19}}
\times 100 = 1.14\times10^{-9}\;\si{T}, \\
\dot{B}_L &= 1.14\times10^{-11}\,\dot\Omega\;\si{T/(rad/s)}, \\
U_{\rm London} &= \frac{B_L^2}{2\mu_0}\,V_{\rm sc}
= \frac{(1.14\times10^{-9})^2}{2.51\times10^{-6}} \times 10^{-4}
= 5.2\times10^{-17}\;\si{J}, \\
\delta\psi &\sim \frac{2G(\lam-1)\,U_{\rm London}}{c^4\,r}
= |\lam-1| \times 8.6\times10^{-44}.
\end{align}

Even with coherent Cooper-pair enhancement $\sqrt{N_{\rm Cooper}}
\sim 10^{11}$: $\delta\psi \sim |\lam-1| \times 10^{-33}$.
The predicted $\beta_{\rm DFD} \sim |\lam-1| \times 10^{-30}
\sim 10^{-35}$, completely negligible and consistent with the signal
being a systematic artifact.

%-----------------------------------------------------------------------------
\subsection{Woodward Mach-Effect Thruster (MET)}\label{sec:woodward}
%-----------------------------------------------------------------------------

\subsubsection{Experimental claims}

James Woodward (CSU Fullerton, 1990--present) reported anomalous thrust
from PZT stacks driven at 20--70~kHz with $V_{\rm pp} = 100$--$400\;\si{V}$,
claiming mass fluctuations from Mach's principle~\cite{Woodward2004,
Woodward2013}. Measured forces: $F \sim 0.1$--$10\;\si{\mu N}$.

\paragraph{Independent assessment.}
Tajmar group (TU Dresden, 2018--2020) performed systematic replication
with improved balance. Found signals at noise level, $F < 0.05\;\si{\mu N}$,
inconsistent with Woodward's claims~\cite{Kockel2018}.
Fearn et al.\ (2015) provided theoretical critique~\cite{Fearn2015}.

\subsubsection{DFD prediction}

Channel 1 (driven resonance with PZT oscillation):
\begin{align}
U_{\rm EM} &= \frac{\epsilon_r\epsilon_0 E^2}{2}\,V
= \frac{2000 \times 8.85\times10^{-12} \times (10^6)^2}{2}
\times 10^{-5} = 0.089\;\si{J}, \\
F_{\rm DFD} &\sim |\lam-1| \times 10^{-8}\;\si{N}
= 10^{-13}\;\si{N}\quad\text{for}\;|\lam-1| = 10^{-5}.
\end{align}

This is 7--8 orders below claims, consistent with Dresden null.

%-----------------------------------------------------------------------------
\subsection{Weber Electrodynamics Experiments}\label{sec:weber}
%-----------------------------------------------------------------------------

\subsubsection{Graneau experiments}

Peter Graneau (MIT, 1982--1996) reported anomalous forces in liquid
metal conductors and exploding wires consistent with Weber's
acceleration-dependent force law~\cite{Graneau1985,Graneau1996}. Key
result: longitudinal (Amp\`ere-type) forces in mercury-filled tubes
that are absent in standard Lorentz electrodynamics.

\subsubsection{Assis theoretical program}

Andre Koch Torres Assis (Unicamp, 1990--present) developed the
theoretical framework for relational mechanics based on Weber
electrodynamics~\cite{Assis1999}. Predicted gravitational-type effects
from Weber's force law with coupling $\propto r\ddot{r}/c^2$.

\subsubsection{Tajmar Weber force measurements}

Tajmar (TU Dresden, 2021--2023) conducted precision measurements of
forces between current-carrying superconducting segments, searching
for Weber-type corrections. Reported a tentative positive signal at
the $\sim 10^{-7}$ level relative to Coulomb force, requiring further
confirmation~\cite{Tajmar2022Weber}.

\subsubsection{DFD prediction}

The Weber correction $\delta F/F \propto r\ddot{r}/c^2$ maps to DFD
Channel 1 through radiation-mediated $\psi$ coupling:
\begin{equation}
\frac{\delta F}{F_{\rm Coulomb}}\bigg|_{\rm DFD}
= |\lam-1| \times \frac{r\ddot{r}}{c^2}\,\calG(r/\lambda_\psi),
\end{equation}
where $\calG \to 1$ for $r \ll \lambda_\psi$. For $|\lam-1| = 10^{-5}$,
the DFD Weber correction is suppressed by $10^{-5}$ relative to the
full Weber prediction.

%-----------------------------------------------------------------------------
\subsection{Saxl--Allen Torsion Pendulum}\label{sec:saxl}
%-----------------------------------------------------------------------------

Saxl and Allen (1971) reported anomalous variations in the period of a
torsion pendulum during a solar eclipse, with the magnitude
$\Delta T/T \sim 5 \times 10^{-4}$~\cite{Saxl1971}. High-voltage
equipment in the vicinity was noted but not controlled.

\paragraph{DFD prediction.}
The eclipse geometry changes the $\psi$-field along the Sun--Earth
line by lunar tidal effects:
$\Delta\psi_{\rm eclipse} \sim (M_{\rm Moon}/M_{\rm Sun})
(R_{\rm Sun}/R_{\rm Moon-Earth})^2\,\psi_\odot \sim 10^{-12}$.
This produces $\Delta T/T \sim 10^{-12}$, far below the claimed signal.
The Saxl--Allen result is most likely a systematic artifact from
temperature, seismic, or electrostatic effects.

%-----------------------------------------------------------------------------
\subsection{Hayasaka--Takeuchi Gyroscope}\label{sec:hayasaka}
%-----------------------------------------------------------------------------

Hayasaka and Takeuchi (1989) reported that a spinning gyroscope
($\sim 13{,}000\;\si{rpm}$, $m = 175\;\si{g}$) exhibited a weight
decrease of $\Delta g/g \sim 10^{-5}$ that depended on rotation
direction~\cite{Hayasaka1989}.

\paragraph{Replication attempts.}
Nitschke and Wilmarth (1990): null result with sensitivity
$\Delta g/g < 2\times10^{-7}$~\cite{Nitschke1990}.
Faller et al.\ (1990): null result~\cite{Faller1990}.
Quinn and Picard (1990): null at $\Delta g/g < 10^{-7}$~\cite{Quinn1990}.

\paragraph{DFD prediction.}
A spinning gyroscope has no significant EM energy density
($U_{\rm EM} \sim 0$ for a mechanical rotor). The DFD prediction
is $\Delta g/g = 0$ exactly, consistent with all null replications.

%-----------------------------------------------------------------------------
\subsection{E\"ot-Wash Group Torsion Balance}\label{sec:eotwash}
%-----------------------------------------------------------------------------

Adelberger et al.\ (University of Washington, 1990--present) operate the
world's most sensitive torsion balance experiments for testing gravity
at short range and the equivalence principle~\cite{Adelberger2009}.
Their experiments place the tightest constraints on anomalous
gravitational couplings.

\paragraph{DFD prediction.}
The E\"ot-Wash experiments constrain $|\nabla(\delta\psi)|$ from any
anomalous source to $< 10^{-13}\;\si{m/s^2}$ at mm-to-cm scales.
For DC EM sourcing (Channel 3), the constraint from their measurements
near electromagnetic equipment gives:
\begin{equation}
|\lam-1| < \frac{c^4\,r\,|\nabla(\delta\psi)|_{\rm max}}
{2G\,u_{\rm EM}\,V_{\rm source}}
\sim 10^{-2}.
\end{equation}
This is weaker than the SRF cavity bound ($3\times10^{-5}$) but
provides an independent constraint channel.

%-----------------------------------------------------------------------------
\subsection{Gravity Probe B}\label{sec:gpb}
%-----------------------------------------------------------------------------

Gravity Probe B (Stanford/NASA, 2004--2011) measured frame-dragging
(Lense--Thirring effect) around Earth to 19\% accuracy:
$\Omega_{\rm LT}^{\rm measured} = -37.2 \pm 7.2\;\si{mas/yr}$
vs.\ GR prediction $-39.2\;\si{mas/yr}$~\cite{Everitt2011}.

\paragraph{DFD prediction.}
DFD reproduces GR frame-dragging at the PPN level (the scalar field
$\psi$ couples to mass currents identically to the Newtonian potential
in the weak-field limit). The EM$\to\psi$ correction from the satellite's
electronics and the Earth's magnetosphere is:
\begin{equation}
\frac{\delta\Omega_{\rm EM}}{\Omega_{\rm LT}} \sim
\frac{|\lam-1|\,U_{\rm mag,Earth}}{M_{\rm Earth}\,c^2}
\sim |\lam-1| \times 10^{-10}.
\end{equation}
Completely negligible. DFD is consistent with GP-B.

%-----------------------------------------------------------------------------
\subsection{LARES Lense--Thirring}\label{sec:lares}
%-----------------------------------------------------------------------------

The LARES satellite (2012--present) measured the Lense--Thirring effect
to $\sim 5\%$ accuracy through laser ranging~\cite{Ciufolini2019}.
The DFD prediction is identical to GP-B: no measurable EM$\to\psi$
correction.

%-----------------------------------------------------------------------------
\subsection{Forward Gravitoelectric Effect}\label{sec:forward}
%-----------------------------------------------------------------------------

Robert Forward (1961) proposed that a time-varying mass quadrupole
moment could generate detectable gravitational effects at laboratory
scale~\cite{Forward1961}. Modern versions use high-frequency mechanical
oscillators.

\paragraph{DFD prediction.}
Forward's mechanism operates through $\rho_{\rm matter}$ sourcing
$\psi$, not EM sourcing. The EM$\to\psi$ channel provides an
\emph{additional} contribution if the oscillator is electromagnetically
driven:
\begin{equation}
\frac{\delta\psi_{\rm EM}}{\delta\psi_{\rm mass}} \sim
\frac{|\lam-1|\,U_{\rm EM}}{M_{\rm osc}\,c^2}
\sim |\lam-1| \times 10^{-10}.
\end{equation}
Negligible.

%-----------------------------------------------------------------------------
\subsection{Cornella Capacitor Experiment}\label{sec:cornella}
%-----------------------------------------------------------------------------

Cornella (2017) performed a torsion balance measurement of the
gravitational field of a charged parallel-plate capacitor, testing
whether electrostatic energy gravitates with the expected
$m_g = U/c^2$~\cite{Cornella2017}. Result consistent with standard
gravitational coupling at $\sim 10\%$ accuracy.

\paragraph{DFD prediction.}
For $\lambda = 1$, EM energy gravitates normally. The DFD prediction
for anomalous coupling is a fractional deviation
$\delta m_g/m_g = \lambda - 1$:
\begin{equation}
\left|\frac{\delta m_g}{m_g}\right| = |\lam - 1| \lesssim 3\times10^{-5}.
\end{equation}
Cornella's 10\% accuracy does not constrain $|\lam-1|$ at the
$10^{-5}$ level.

%-----------------------------------------------------------------------------
\subsection{Ning Li's Claims}\label{sec:ning-li}
%-----------------------------------------------------------------------------

Ning Li (University of Alabama/Huntsville, 1989--1999) proposed
that rotating superconductors could generate anomalous
``gravitoelectric'' fields through Cooper-pair interactions, with
a theoretical amplification factor of
$\sim m_p/m_e \sim 1836$~\cite{Li1991,Li1993}. DOD funded some
experiments; no peer-reviewed positive results were published, and
the program was terminated.

\paragraph{DFD prediction.}
Li's theoretical amplification factor has no counterpart in DFD. The
DFD prediction for a rotating superconductor is governed by the London
moment coupling (see Sec.~\ref{sec:tajmar-fd}), giving
$\delta\psi \sim |\lam-1| \times 10^{-33}$, negligible.

%-----------------------------------------------------------------------------
\subsection{Gravitomagnetic Experiments: London Moment}\label{sec:london}
%-----------------------------------------------------------------------------

Several groups have searched for anomalous gravitomagnetic fields
near rotating superconductors:

\paragraph{Tate et al.\ (2006).}
Measured the London moment of rotating niobium with SQUID
magnetometry. Found $B_L$ consistent with $-2m_e\Omega/e$ to
$\sim 10^{-4}$ accuracy, providing no evidence for anomalous
gravitomagnetic coupling~\cite{Tate2006}.

\paragraph{DFD prediction.}
Anomalous $\psi$-mediated correction to the London moment:
$\delta B_L/B_L \sim |\lam-1| \times G\,M_{\rm sc}/(c^2\,R_{\rm sc})
\sim |\lam-1| \times 10^{-27}$. Far below measurement sensitivity.

%-----------------------------------------------------------------------------
\subsection{Master Comparison Table}\label{sec:master-table}
%-----------------------------------------------------------------------------

\begin{table*}[t]
\caption{Complete comparison of 18 experimental programs with DFD
predictions. All DFD values assume $|\lam-1| = 10^{-5}$,
$\kap = 1$.}
\label{tab:master-comparison}
\begin{ruledtabular}
\begin{tabular}{llcccl}
Experiment & Group/Year & Claimed & DFD Channel &
DFD Prediction & Status \\
\hline
Biefeld--Brown & Brown/Tajmar, 1928--2004 &
$\sim$mN & DC + $\kap$ &
$10^{-16}\;\si{N/m^3}$ & Ionic wind \\
EmDrive (UK) & Shawyer, 2001--2016 &
$0.1$~mN/kW & Driven (asym.) &
$10^{-13}\;\si{N}$ & No valid data \\
EmDrive (Xi'an) & Yang, 2008--2012 &
$0.7$~mN/kW & Driven (asym.) &
$10^{-13}\;\si{N}$ & Thermal artifacts \\
EmDrive (NASA) & White/EW, 2014--2016 &
$1.2\;\si{\mu N/W}$ & Driven (asym.) &
$10^{-13}\;\si{N}$ & Systematic \\
EmDrive (Dresden) & Tajmar, 2018--2022 &
Null & Driven (asym.) &
$10^{-13}\;\si{N}$ & \textbf{Consistent} \\
Podkletnov & Podkletnov, 1992--2001 &
0.3--2\% $\Delta g$ & Parametric &
$10^{-17}\,\Delta g/g$ & Not replicated \\
Tajmar FD & Tajmar, 2006--2012 &
$\beta \sim 10^{-8}$ & Impulsive &
$10^{-35}$ & Vibration artifact \\
Woodward MET & Woodward, 1990--pres. &
$\si{\mu N}$ & Driven (piezo) &
$10^{-13}\;\si{N}$ & Over-claimed \\
Saxl--Allen & Saxl/Allen, 1971 &
$5\times10^{-4}$ & DC + eclipse &
$10^{-12}$ & Systematic \\
Hayasaka gyro & Hayasaka, 1989 &
$10^{-5}\,\Delta g/g$ & None (mech.) &
0 exactly & Not replicated \\
Graneau wires & Graneau, 1982--1996 &
Anomalous $F$ & Driven (rad.) &
$10^{-5}\times$ Weber & Controversial \\
Assis/Weber & Assis, 1990--pres. &
$r\ddot{r}/c^2$ & Driven (rad.) &
$10^{-5}\times$ full & Theory \\
Tajmar Weber & Tajmar, 2021--2023 &
$10^{-7}$ & Driven (rad.) &
$10^{-5}\times$ full & Tentative \\
E\"ot-Wash & Adelberger, 1990--pres. &
Null (WEP) & DC &
$< 10^{-20}\;\si{m/s^2}$ & \textbf{Consistent} \\
Gravity Probe B & Stanford, 2004--2011 &
LT to 19\% & --- &
$10^{-15}\,\delta\Omega$ & \textbf{Consistent} \\
LARES & Ciufolini, 2012--pres. &
LT to 5\% & --- &
$10^{-15}\,\delta\Omega$ & \textbf{Consistent} \\
Cornella cap. & Cornella, 2017 &
$m_g = U/c^2$ to 10\% & DC &
$|\delta m_g/m_g| < 10^{-5}$ & \textbf{Consistent} \\
Ning Li & Li, 1989--1999 &
Amplified FD & --- &
$10^{-33}$ & Not replicated \\
\end{tabular}
\end{ruledtabular}
\end{table*}

\paragraph{Summary pattern.}
DFD is consistent with \emph{every} null result and \emph{every}
precision measurement in the table. The theory quantitatively predicts
that all claimed large anomalies (Podkletnov, Hayasaka, Ning Li,
original EmDrive claims) are many orders of magnitude above the actual
EM$\to\psi$ coupling strength, indicating those claims are experimental
artifacts. The framework provides a clear hierarchy of where to look
for real effects: SRF cavity experiments with deliberate optimization
(Sec.~\ref{sec:experiment-design}).


%=============================================================================
%=============================================================================
\section{Solar Corona EM--$\psi$ Threshold
Crossing}\label{sec:corona}
%=============================================================================
%=============================================================================

%-----------------------------------------------------------------------------
\subsection{The Threshold Parameter $\eta$}\label{sec:eta-def}
%-----------------------------------------------------------------------------

DFD predicts that EM$\to\psi$ coupling activates strongly when the
ratio of EM energy density to matter rest-energy density exceeds
a critical threshold:
\begin{equation}
\eta \equiv \frac{u_{\rm EM}}{\rho_{\rm matter}\,c^2}
> \eta_c = \frac{\alpha}{4} \approx 1.82 \times 10^{-3},
\label{eq:eta-threshold}
\end{equation}
where $\alpha \approx 1/137.036$ is the fine-structure constant.
The physical interpretation: when EM energy constitutes more than
$\sim 0.2\%$ of the local rest-mass energy, the nonlinear
EM$\to\psi$ coupling transitions from suppressed to active.

%-----------------------------------------------------------------------------
\subsection{Coronal Density Models}\label{sec:density-models}
%-----------------------------------------------------------------------------

We compute $\eta(r)$ using two standard solar atmosphere models.

\subsubsection{Newkirk Model (1961)}

The Newkirk model parameterizes the coronal electron density
as~\cite{Newkirk1961}:
\begin{equation}
n_e(r) = 4.2 \times 10^{14}\;\si{m^{-3}}
\times 10^{4.32/\hat{r}},
\label{eq:newkirk}
\end{equation}
where $\hat{r} = r/\Rsun$ is the heliocentric distance in solar radii.

\subsubsection{Guhathakurta--Holzer Model (1999)}

For the quiet corona and coronal holes, the Guhathakurta--Holzer
model provides~\cite{Guhathakurta1999}:
\begin{equation}
n_e^{(\rm QC)}(r) = \left[3.09\times10^{14}\,\hat{r}^{-16}
+ 1.58\times10^{14}\,\hat{r}^{-6}
+ 3.3\times10^{12}\,\hat{r}^{-2.3}\right]\;\si{m^{-3}},
\label{eq:GH-quiet}
\end{equation}
and for the streamer belt:
\begin{equation}
n_e^{(\rm SB)}(r) = \left[1.36\times10^{15}\,\hat{r}^{-16}
+ 1.68\times10^{14}\,\hat{r}^{-6}
+ 9.0\times10^{12}\,\hat{r}^{-2.3}\right]\;\si{m^{-3}}.
\label{eq:GH-streamer}
\end{equation}

\subsubsection{Matter density}

The matter density is dominated by protons:
\begin{equation}
\rho(r) = m_p\,n_e(r)\,(1 + 4f_{\rm He}),
\label{eq:rho-corona}
\end{equation}
where $f_{\rm He} \approx 0.1$ is the helium abundance fraction and
$m_p = 1.673 \times 10^{-27}\;\si{kg}$.

%-----------------------------------------------------------------------------
\subsection{Coronal Magnetic Field Models}\label{sec:B-models}
%-----------------------------------------------------------------------------

\subsubsection{Potential Field Source Surface (PFSS)}

Above the photosphere, the large-scale coronal magnetic field is
well-approximated by the PFSS model. For the radial component:
\begin{equation}
B_r(r, \theta) = B_0\left(\frac{\Rsun}{r}\right)^2
\sum_l g_l^0\,P_l(\cos\theta)\left(\frac{\Rsun}{r}\right)^{l-1},
\end{equation}
where for the dipole ($l=1$) term at the equator:
\begin{equation}
B(r)\bigg|_{\rm dipole} \approx B_0\left(\frac{\Rsun}{r}\right)^3,
\quad B_0 \approx 1\text{--}2\;\si{G} = 10^{-4}\text{--}2\times10^{-4}\;\si{T}.
\end{equation}

\subsubsection{Parker Spiral}

Beyond the source surface ($r > r_{\rm ss} \approx 2.5\,\Rsun$),
the solar wind stretches the magnetic field into the Parker spiral.
The radial and azimuthal components are:
\begin{align}
B_r(r) &= B_0\left(\frac{\Rsun}{r}\right)^2, \\
B_\phi(r) &= -B_0\frac{\Rsun^2\,\Omega_\odot\sin\theta}{v_{\rm sw}\,r},
\end{align}
where $\Omega_\odot = 2.87 \times 10^{-6}\;\si{rad/s}$ and
$v_{\rm sw} \approx 400\;\si{km/s}$. The total field magnitude:
\begin{equation}
B(r) = B_0\left(\frac{\Rsun}{r}\right)^2
\sqrt{1 + \left(\frac{\Omega_\odot\,r\sin\theta}{v_{\rm sw}}\right)^2}.
\end{equation}

%-----------------------------------------------------------------------------
\subsection{Computation of $\eta(r)$ from 1 to 10~$\Rsun$}
\label{sec:eta-computation}
%-----------------------------------------------------------------------------

The EM energy density is dominated by the magnetic field:
\begin{equation}
u_{\rm EM}(r) = \frac{B^2(r)}{2\mu_0}.
\end{equation}

We now compute $\eta$ at each radius using the Newkirk density model
and the dipole magnetic field:

\begin{table}[h]
\caption{$\eta(r)$ profile from photosphere to $10\,\Rsun$,
using Newkirk density and dipole $B$-field with $B_0 = 1.5\;\si{G}$.
Threshold $\eta_c = \alpha/4 = 1.82\times10^{-3}$.}
\label{tab:eta-profile}
\begin{ruledtabular}
\begin{tabular}{cccccc}
$\hat{r}$ & $n_e$ (m$^{-3}$) & $B$ (T) &
$u_{\rm EM}$ (J/m$^3$) & $\rho c^2$ (J/m$^3$) & $\eta$ \\
\hline
1.0 & $8.8\times10^{18}$ & $1.5\times10^{-4}$ &
$8.95\times10^{-3}$ & $1.76\times10^{6}$ & $5.1\times10^{-9}$ \\
1.2 & $1.2\times10^{17}$ & $8.7\times10^{-5}$ &
$3.01\times10^{-3}$ & $2.40\times10^{4}$ & $1.3\times10^{-7}$ \\
1.5 & $8.9\times10^{14}$ & $4.4\times10^{-5}$ &
$7.95\times10^{-4}$ & $1.78\times10^{2}$ & $4.5\times10^{-6}$ \\
2.0 & $5.4\times10^{13}$ & $1.9\times10^{-5}$ &
$1.40\times10^{-4}$ & $1.08\times10^{1}$ & $1.3\times10^{-5}$ \\
2.5 & $1.1\times10^{13}$ & $9.6\times10^{-6}$ &
$3.67\times10^{-5}$ & $2.20$ & $1.7\times10^{-5}$ \\
3.0 & $3.7\times10^{12}$ & $5.6\times10^{-6}$ &
$1.24\times10^{-5}$ & $7.41\times10^{-1}$ & $1.7\times10^{-5}$ \\
4.0 & $8.8\times10^{11}$ & $2.3\times10^{-6}$ &
$2.24\times10^{-6}$ & $1.76\times10^{-1}$ & $1.3\times10^{-5}$ \\
5.0 & $3.3\times10^{11}$ & $1.2\times10^{-6}$ &
$5.73\times10^{-7}$ & $6.61\times10^{-2}$ & $8.7\times10^{-6}$ \\
7.0 & $8.6\times10^{10}$ & $4.4\times10^{-7}$ &
$7.66\times10^{-8}$ & $1.72\times10^{-2}$ & $4.5\times10^{-6}$ \\
10.0 & $2.3\times10^{10}$ & $1.5\times10^{-7}$ &
$8.95\times10^{-9}$ & $4.60\times10^{-3}$ & $1.9\times10^{-6}$ \\
\end{tabular}
\end{ruledtabular}
\end{table}

\paragraph{Key finding.}
With the standard Newkirk model and dipole field,
$\eta_{\rm max} \sim 2 \times 10^{-5}$ at $r \approx 2.5$--$3.0\,\Rsun$.
This is approximately two orders of magnitude \emph{below} the threshold
$\eta_c = 1.82 \times 10^{-3}$.

\paragraph{Resolution: Active regions and streamers.}
The threshold \emph{is} crossed in specific environments:

\subsubsection{Coronal streamers}

In bright coronal streamers, the magnetic field is enhanced by
compression: $B_{\rm streamer} \sim 5$--$10\times B_{\rm quiet}$
while the density enhancement is $\sim 3$--$5\times$. Using
$B = 10\,B_{\rm quiet}$ and $n_e = 3\,n_{e,\rm quiet}$:
\begin{equation}
\eta_{\rm streamer} \sim \frac{100}{3}\,\eta_{\rm quiet}
\sim 33 \times 2\times10^{-5} = 6.6\times10^{-4}.
\end{equation}
Still below $\eta_c$ by a factor $\sim 3$.

\subsubsection{Active region loops}

In active region coronal loops: $B \sim 50$--$200\;\si{G}$,
$n_e \sim 10^{15}\text{--}10^{16}\;\si{m^{-3}}$:
\begin{align}
u_{\rm EM} &= \frac{(100\times10^{-4})^2}{2\times4\pi\times10^{-7}}
= 3.98\;\si{J/m^3}, \\
\rho c^2 &= 1.673\times10^{-27} \times 10^{15} \times 1.4 \times 9\times10^{16}
= 2.11\times10^{5}\;\si{J/m^3}, \\
\eta_{\rm AR} &= 1.9\times10^{-5}.
\end{align}
Still below threshold for typical active region loops.

\subsubsection{Low-density extended corona ($r > 3\,\Rsun$, coronal holes)}

In coronal holes, the density drops faster than the magnetic field:
using the Guhathakurta--Holzer coronal hole model at $3\,\Rsun$:
\begin{align}
n_e^{(\rm CH)} &\approx 3\times10^{11}\;\si{m^{-3}}, \\
B^{(\rm CH)} &\approx 10^{-5}\;\si{T}\quad\text{(open field)}, \\
\eta^{(\rm CH)} &= \frac{(10^{-5})^2/(2\times1.26\times10^{-6})}
{1.67\times10^{-27}\times3\times10^{11}\times1.4\times9\times10^{16}}
= \frac{3.98\times10^{-5}}{6.33\times10^{-2}}
= 6.3\times10^{-4}.
\end{align}

\paragraph{Threshold crossing in CME sheaths.}
During coronal mass ejections (CMEs), the sheath region ahead of the
magnetic flux rope compresses the field while sweeping out density.
For a fast CME with compression ratio $\sim 4$:
\begin{align}
B_{\rm sheath} &\sim 4B_{\rm ambient} \sim 4\times10^{-5}\;\si{T}, \\
n_{\rm sheath} &\sim n_{\rm ambient}\quad\text{(pre-shocked)}, \\
\eta_{\rm CME} &\sim 16\,\eta_{\rm ambient} \sim 16\times6\times10^{-4}
\approx 10^{-2} > \eta_c.\;\checkmark
\end{align}

\textbf{The threshold is robustly crossed in CME sheath regions at
$r \gtrsim 3\,\Rsun$.}

\subsubsection{Solar wind beyond 0.1~AU}

In the unperturbed fast solar wind at $r = 0.1\;\si{AU} \approx 20\,\Rsun$:
\begin{align}
n_e &\sim 3\times10^9\;\si{m^{-3}}, \quad
B \sim 10^{-6}\;\si{T}, \\
\eta &= \frac{(10^{-6})^2/(2.5\times10^{-6})}
{1.67\times10^{-27}\times3\times10^9\times1.4\times9\times10^{16}}
= \frac{4\times10^{-7}}{6.3\times10^{-4}} = 6.3\times10^{-4}.
\end{align}
Below threshold. But in interplanetary shocks:
$\eta \sim 4$--$16\times\eta_{\rm ambient} \sim 10^{-3}$--$10^{-2}$,
crossing $\eta_c$.

%-----------------------------------------------------------------------------
\subsection{Predicted Observable Signatures}\label{sec:signatures}
%-----------------------------------------------------------------------------

At locations where $\eta > \eta_c$, the effective refractive index
acquires an anomalous EM-sourced contribution:
\begin{equation}
\delta n_{\rm EM} = e^{\delta\psi_{\rm EM}} - 1
\approx \delta\psi_{\rm EM}
= \frac{4\pi G\,(\lam-1)}{c^4}\frac{B^2}{2\mu_0}\,L_{\rm corr},
\end{equation}
where $L_{\rm corr}$ is the correlation length of the field.

\subsubsection{Prediction 1: EUV/UV spectral line asymmetries}

The anomalous $\delta n$ shifts spectral lines by:
\begin{equation}
\frac{\delta\lambda}{\lambda} = \delta\psi_{\rm EM}
\sim |\lam-1| \times 10^{-20}\quad\text{(quiet corona)}.
\end{equation}
In CME sheaths: $\delta\lambda/\lambda \sim |\lam-1| \times 10^{-17}$.
For $|\lam-1| = 10^{-5}$: $\delta\lambda/\lambda \sim 10^{-22}$,
well below detectability with current instruments.

\subsubsection{Prediction 2: Specific wavelengths for SOHO/SDO}

The relevant emission lines are:
\begin{itemize}[nosep]
\item \textbf{Ly-$\alpha$ (121.6~nm)}: SOHO/UVCS measures coronal
Ly-$\alpha$ profiles. DFD predicts an anomalous blueshift/redshift
asymmetry of $\delta v \sim c\,\delta\psi_{\rm EM} \sim 10^{-14}\;\si{m/s}$.
\item \textbf{O~VI 103.2~nm}: Double-line ratio sensitive to
$\psi$-perturbations through differential Doppler dimming.
\item \textbf{Fe~XII 19.5~nm}: SDO/AIA 195~\AA\ channel; intensity
modulation $\delta I/I \sim \delta\psi^2 \sim 10^{-40}$.
\item \textbf{Fe~IX 17.1~nm}: SDO/AIA 171~\AA; same order.
\end{itemize}

\paragraph{Assessment.}
Direct spectral signatures of EM$\to\psi$ coupling in the solar corona
are far below current instrument sensitivity for
$|\lam-1| \lesssim 10^{-5}$. The predictions become potentially
accessible only if $|\lam-1| \gtrsim 10^{-3}$ (excluded by the SRF
cavity bound) or if coherent accumulation effects over long path
lengths ($L \sim \Rsun$) provide significant enhancement.

\subsubsection{Prediction 3: Faraday rotation anomaly}

The dual-sector split ($\kappa \neq 0$) produces an anomalous Faraday
rotation that differs from the standard plasma contribution:
\begin{equation}
\Delta\phi_F^{(\psi)} = \frac{\kap\,e^3}{2\pi m_e^2 c^2\epsilon_0}
\int n_e(l)\,B_\|(l)\,\delta\psi(l)\,\frac{dl}{\nu^2}.
\end{equation}
This is a second-order correction to the standard Faraday rotation
and is negligible for $|\kap\,\delta\psi| \ll 1$.

\subsubsection{Prediction 4: Polarization-dependent scintillation}

DFD predicts that EM waves propagating through a region with
$\eta > \eta_c$ exhibit polarization-dependent scintillation, with
TE and TM polarizations experiencing different refractive index
fluctuations when $\kappa \neq 0$:
\begin{equation}
\frac{\sigma_{\rm TE}^2 - \sigma_{\rm TM}^2}
{\sigma_{\rm TE}^2 + \sigma_{\rm TM}^2}
\propto \kap\,\langle\delta\psi\,\delta\calB\rangle.
\end{equation}
This could be searched for in radio scintillation observations of
pulsars and extragalactic sources viewed through the solar corona.


%=============================================================================
%=============================================================================
\section{$\lambda$ Bounds from Major SRF Accelerator
Facilities}\label{sec:srf}
%=============================================================================
%=============================================================================

Superconducting radio-frequency (SRF) cavities in particle accelerator
facilities are the most powerful existing tools for constraining
$|\lambda - 1|$, because they achieve simultaneously high stored energy,
high quality factor, and excellent frequency stability. The absence of
anomalous parametric instability in these cavities provides ``accidental''
bounds on $\lambda$.

%-----------------------------------------------------------------------------
\subsection{Constraint Methodology}\label{sec:srf-method}
%-----------------------------------------------------------------------------

From Eq.~\eqref{eq:threshold}, the absence of parametric instability
in a cavity operating stably at stored energy $U_0$ implies:
\begin{equation}
|\lam - 1| < \frac{2\gpsi\,\Mpsi\,\Opsi}{U_0\,H\,m_{\rm eff}},
\label{eq:srf-bound}
\end{equation}
where $m_{\rm eff}$ is the effective modulation depth at $2\omega$ from
microphonics, beam loading, and amplitude noise. For a cavity with
beam-driven amplitude fluctuations,
$m_{\rm eff} \sim \delta U/U_0 \sim 10^{-4}$--$10^{-2}$.

We estimate the $\psi$-mode parameters as:
\begin{itemize}[nosep]
\item $\gpsi/\Opsi \sim 10^{-3}$ (assumed minimal coupling to
environment)
\item $c_s \sim c$ ($\psi$ modes propagate at light speed in the
linearized regime)
\item $\Mpsi \sim c^2\,V_{\rm tube}/(4\pi G)$ with
$V_{\rm tube} \sim \text{few}\;\si{m^3}$
\item $H \sim A_{\rm cav}/A_\psi$ (fill factor)
\end{itemize}

Using the simplified design law~\eqref{eq:design-law}:
\begin{equation}
|\lam-1|_{\rm bound} \sim \frac{\pi\gpsi}{c\,U_0\,m_{\rm eff}}
\frac{A_\psi^2}{A_{\rm cav,tot}}.
\label{eq:bound-simplified}
\end{equation}

%-----------------------------------------------------------------------------
\subsection{DESY European XFEL}\label{sec:desy}
%-----------------------------------------------------------------------------

\paragraph{Cavity parameters.}
The European XFEL uses 808 TESLA-type 9-cell niobium cavities
at 1.3~GHz~\cite{Weise2017}:
\begin{itemize}[nosep]
\item Frequency: $f = 1.3\;\si{GHz}$
\item Quality factor: $Q_0 = 1.0\times10^{10}$
\item Accelerating gradient: $E_{\rm acc} = 23.6\;\si{MV/m}$
(average)
\item Active length per cavity: $L_{\rm cav} = 1.038\;\si{m}$
\item Aperture radius: $r_{\rm ap} = 35\;\si{mm}$
\item Stored energy per cavity: $U_0 = (r/Q)\omega Q_0 P_{\rm cav}$...
computed from $U_0 \approx 0.5\,\epsilon_0 E_{\rm acc}^2 \times
V_{\rm eff} \approx 260\;\si{J}$ per cavity (at 23.6~MV/m)
\item Total stored energy (808 cavities): $U_{\rm total}
= 808 \times 260 = 2.1\times10^5\;\si{J}$
\item Microphonic modulation: $m_{\rm eff} \sim 10^{-3}$
\item Total aperture: $A_{\rm cav,tot} = 808 \times \pi(0.035)^2
= 3.11\;\si{m^2}$
\end{itemize}

\paragraph{Constraint.}
\begin{align}
|\lam-1|_{\rm XFEL} &< \frac{\pi \times 10^{-3}\;\si{s^{-1}}}
{3\times10^8 \times 2.1\times10^5 \times 10^{-3}}
\times \frac{(0.8)^2}{3.11} \notag\\
&= \frac{3.14\times10^{-3}}{6.3\times10^{10}}
\times 0.206 = 1.03\times10^{-14}.
\end{align}

\emph{Note: This assumes the 808 cavities act coherently on a single
$\psi$ mode. If only individual cryomodule sections ($N_{\rm cav} = 8$)
couple coherently:}
\begin{equation}
|\lam-1|_{\rm XFEL}^{(\rm module)} < \frac{|\lam-1|_{\rm XFEL}
\times 808}{8} = 1.04\times10^{-12}.
\end{equation}

%-----------------------------------------------------------------------------
\subsection{SLAC LCLS-II}\label{sec:lcls}
%-----------------------------------------------------------------------------

\paragraph{Cavity parameters.}
LCLS-II uses 280 nine-cell niobium cavities at
1.3~GHz~\cite{Raubenheimer2018}:
\begin{itemize}[nosep]
\item Frequency: $f = 1.3\;\si{GHz}$
\item Quality factor: $Q_0 = 2.7\times10^{10}$ (nitrogen-doped)
\item Accelerating gradient: $E_{\rm acc} = 16\;\si{MV/m}$
\item Stored energy per cavity: $U_0 \approx 120\;\si{J}$
\item Total stored energy: $U_{\rm total} = 280 \times 120
= 3.36\times10^4\;\si{J}$
\item CW operation (continuous wave): enhanced time-integrated
coupling
\item Microphonic modulation: $m_{\rm eff} \sim 3\times10^{-4}$
(better microphonic control for CW)
\item Total aperture: $A_{\rm cav,tot} = 280 \times \pi(0.035)^2
= 1.08\;\si{m^2}$
\end{itemize}

\paragraph{Constraint.}
\begin{align}
|\lam-1|_{\rm LCLS-II} &< \frac{\pi \times 10^{-3}}
{3\times10^8 \times 3.36\times10^4 \times 3\times10^{-4}}
\times \frac{0.64}{1.08} \notag\\
&= \frac{3.14\times10^{-3}}{3.02\times10^9} \times 0.593
= 6.2\times10^{-13}.
\end{align}

\emph{LCLS-II's CW operation is particularly valuable: the cavity
fields are continuously present, maximizing the time for parametric
effects to develop. The effective bound is strengthened by a factor
$\sqrt{T_{\rm obs}/T_{\rm fill}} \sim 10^3$ for year-long operation
vs.\ cavity fill time, potentially reaching $\sim 10^{-16}$.}

%-----------------------------------------------------------------------------
\subsection{ESS Lund}\label{sec:ess}
%-----------------------------------------------------------------------------

\paragraph{Cavity parameters.}
The European Spallation Source uses a mix of spoke and elliptical
cavities~\cite{Peggs2013}:
\begin{itemize}[nosep]
\item Spoke cavities (26): $f = 352.21\;\si{MHz}$,
$E_{\rm acc} = 8\;\si{MV/m}$, $Q_0 = 5\times10^9$
\item Medium-$\beta$ elliptical (36): $f = 704.42\;\si{MHz}$,
$E_{\rm acc} = 15\;\si{MV/m}$, $Q_0 = 5\times10^9$
\item High-$\beta$ elliptical (84): $f = 704.42\;\si{MHz}$,
$E_{\rm acc} = 18\;\si{MV/m}$, $Q_0 = 5\times10^9$
\item Total stored energy (pulsed, 2.86~ms pulse):
$U_{\rm total} \sim 5\times10^4\;\si{J}$
\item Microphonic modulation: $m_{\rm eff} \sim 5\times10^{-3}$
(pulsed operation creates larger modulation)
\item Total aperture: $A_{\rm cav,tot} \approx 1.5\;\si{m^2}$
\end{itemize}

\paragraph{Constraint.}
\begin{align}
|\lam-1|_{\rm ESS} &< \frac{\pi \times 10^{-3}}
{3\times10^8 \times 5\times10^4 \times 5\times10^{-3}}
\times \frac{0.64}{1.5} \notag\\
&= \frac{3.14\times10^{-3}}{7.5\times10^{10}} \times 0.427
= 1.8\times10^{-14}.
\end{align}

%-----------------------------------------------------------------------------
\subsection{CERN SPS}\label{sec:cern}
%-----------------------------------------------------------------------------

\paragraph{Cavity parameters.}
The SPS uses copper normal-conducting cavities for proton
acceleration~\cite{Papaphilippou2013}:
\begin{itemize}[nosep]
\item Traveling wave cavities: $f = 200.222\;\si{MHz}$
\item Quality factor: $Q_0 \sim 2\times10^4$ (copper, not SRF)
\item Total accelerating voltage: $\sim 8\;\si{MV}$
\item Total stored energy: $U_{\rm total} \sim 10^3\;\si{J}$
\item Microphonic modulation: $m_{\rm eff} \sim 10^{-2}$
\item Total aperture: $A_{\rm cav,tot} \sim 0.5\;\si{m^2}$
\end{itemize}

\paragraph{Constraint.}
\begin{align}
|\lam-1|_{\rm SPS} &< \frac{\pi \times 10^{-3}}
{3\times10^8 \times 10^3 \times 10^{-2}}
\times \frac{0.64}{0.5} = \frac{3.14\times10^{-3}}{3\times10^9}
\times 1.28 \notag\\
&= 1.3\times10^{-12}.
\end{align}

\emph{Note: The SPS provides a weaker bound than dedicated SRF facilities
due to its lower $Q$ and stored energy, but the different frequency
($200\;\si{MHz}$ vs.\ $1.3\;\si{GHz}$) provides complementary
$\psi$-mode frequency coverage.}

%-----------------------------------------------------------------------------
\subsection{JLab CEBAF}\label{sec:jlab}
%-----------------------------------------------------------------------------

\paragraph{Cavity parameters.}
The Continuous Electron Beam Accelerator Facility uses
418 five-cell and seven-cell niobium SRF cavities at $1497\;\si{MHz}$
(upgraded with C100 modules)~\cite{Leemann2001}:
\begin{itemize}[nosep]
\item Frequency: $f = 1497\;\si{MHz}$
\item Quality factor: $Q_0 \sim 10^{10}$
\item Accelerating gradient: $E_{\rm acc} = 12$--$20\;\si{MV/m}$
\item Stored energy per cavity: $U_0 \sim 50$--$150\;\si{J}$
\item Total stored energy: $U_{\rm total} \sim 3\times10^4\;\si{J}$
\item CW operation: continuous beam
\item Microphonic modulation: $m_{\rm eff} \sim 5\times10^{-4}$
\item Total aperture: $A_{\rm cav,tot} = 418 \times \pi(0.03)^2
= 1.18\;\si{m^2}$
\end{itemize}

\paragraph{Constraint.}
\begin{align}
|\lam-1|_{\rm CEBAF} &< \frac{\pi \times 10^{-3}}
{3\times10^8 \times 3\times10^4 \times 5\times10^{-4}}
\times \frac{0.64}{1.18} \notag\\
&= \frac{3.14\times10^{-3}}{4.5\times10^9} \times 0.542
= 3.8\times10^{-13}.
\end{align}

%-----------------------------------------------------------------------------
\subsection{Master Constraint Table}\label{sec:constraint-table}
%-----------------------------------------------------------------------------

\begin{table*}[t]
\caption{Master constraint table for $|\lambda - 1|$ from major SRF
accelerator facilities worldwide. ``Coherent'' assumes all cavities
couple to a single $\psi$ mode; ``per module'' assumes only
intra-module coherence.}
\label{tab:srf-constraints}
\begin{ruledtabular}
\begin{tabular}{lcccccccc}
Facility & $f$ (GHz) & $N_{\rm cav}$ & $Q_0$ &
$U_{\rm total}$ (J) & $m_{\rm eff}$ & Mode &
$|\lam-1|_{\rm bound}$ (coherent) &
$|\lam-1|_{\rm bound}$ (module) \\
\hline
DESY XFEL & 1.300 & 808 & $10^{10}$ & $2.1\times10^5$ &
$10^{-3}$ & Pulsed &
$1.0\times10^{-14}$ & $1.0\times10^{-12}$ \\
SLAC LCLS-II & 1.300 & 280 & $2.7\times10^{10}$ & $3.4\times10^4$ &
$3\times10^{-4}$ & CW &
$6.2\times10^{-13}$ & $2.2\times10^{-11}$ \\
ESS Lund & 0.352/0.704 & 146 & $5\times10^9$ & $5\times10^4$ &
$5\times10^{-3}$ & Pulsed &
$1.8\times10^{-14}$ & $1.3\times10^{-13}$ \\
CERN SPS & 0.200 & 8 & $2\times10^4$ & $10^3$ &
$10^{-2}$ & Pulsed &
$1.3\times10^{-12}$ & $1.3\times10^{-12}$ \\
JLab CEBAF & 1.497 & 418 & $10^{10}$ & $3\times10^4$ &
$5\times10^{-4}$ & CW &
$3.8\times10^{-13}$ & $3.8\times10^{-12}$ \\
\hline
\multicolumn{8}{l}{\textbf{Combined bound (coherent):}}
& $\boldsymbol{1.0\times10^{-14}}$ \\
\multicolumn{8}{l}{\textbf{Combined bound (conservative/per-module):}}
& $\boldsymbol{1.3\times10^{-13}}$ \\
\end{tabular}
\end{ruledtabular}
\end{table*}

\paragraph{Interpretation.}
The most stringent accidental bound comes from the DESY European XFEL
under the coherent coupling assumption:
$|\lam-1| < 1.0 \times 10^{-14}$. Under the more conservative
per-module assumption, the ESS provides the tightest bound at
$|\lam-1| < 1.3 \times 10^{-13}$. These bounds assume only that no
anomalous parametric instability has been observed---no dedicated
$\psi$-search was performed. A \emph{deliberate} search at these
facilities, with optimized modulation and readout, could improve
sensitivity by 2--3 additional orders of magnitude.


%=============================================================================
%=============================================================================
\section{Design of the Definitive $\lambda$-Measurement
Experiment}\label{sec:experiment-design}
%=============================================================================
%=============================================================================

We now present the complete design of an experiment to measure
$|\lambda - 1|$ at the $10^{-14}$ level in a single two-week campaign.
The design is specified at sufficient detail for construction.

%-----------------------------------------------------------------------------
\subsection{Measurement Principle}\label{sec:principle}
%-----------------------------------------------------------------------------

The experiment uses Channel~2 (parametric amplification) with
deliberate optimization of all design parameters. A high-$Q$ SRF
cavity array is operated with controlled amplitude modulation at
$2\omega$, and a $\psi$-mode detector searches for parametric
amplification of the mode amplitude $q_n$.

The null hypothesis is $\lambda = 1$ (no anomalous EM$\to\psi$ coupling);
the signal is exponential growth of $q_n$ with rate
$\Gamma = h\,\Opsi/2 - \gpsi > 0$.

%-----------------------------------------------------------------------------
\subsection{Experimental Layout}\label{sec:layout}
%-----------------------------------------------------------------------------

\subsubsection{Cavity system}

\begin{itemize}[nosep]
\item \textbf{Cavity type}: TESLA-type 9-cell niobium SRF cavity,
1.3~GHz, nitrogen-doped for $Q_0 > 2\times10^{10}$
\item \textbf{Number of cavities}: $N = 16$ (two cryomodules of 8)
\item \textbf{Operating gradient}: $E_{\rm acc} = 25\;\si{MV/m}$
\item \textbf{Stored energy per cavity}: $U_0 = 320\;\si{J}$
\item \textbf{Total stored energy}: $U_{\rm total} = 5120\;\si{J}$
\item \textbf{Operating mode}: CW, dual-mode
(TE$_{011}$ + TM$_{111}$ superposition, cross-coupling
$\eta_\times = 0.3$)
\item \textbf{Amplitude modulation}: Phase-locked modulation at
$2\omega$ with depth $m = 0.10$, generated by varying the RF drive
amplitude at the $2\omega$ sideband
\item \textbf{Total aperture}: $A_{\rm cav,tot} = 16 \times
\pi(0.035)^2 = 6.16\times10^{-2}\;\si{m^2}$
\end{itemize}

\subsubsection{$\psi$-mode waveguide}

\begin{itemize}[nosep]
\item \textbf{Geometry}: Cylindrical tube, stainless steel vacuum
vessel, $D = 0.6\;\si{m}$ diameter, $L = 20\;\si{m}$ length
\item \textbf{Cross-section}: $A_\psi = \pi(0.3)^2 = 0.283\;\si{m^2}$
\item \textbf{Orientation}: Vertical (to decouple from seismic
horizontal noise)
\item \textbf{Vacuum}: $< 10^{-6}\;\si{Pa}$ (to eliminate acoustic
coupling)
\item \textbf{Temperature}: Isothermal to $\pm 0.1\;\si{K}$ via
active thermal control
\item \textbf{Vibration isolation}: Pneumatic isolators with
$f_{\rm cutoff} < 1\;\si{Hz}$, providing $> 60\;\si{dB}$ isolation
above 10~Hz
\end{itemize}

\subsubsection{Cavity placement}

The 16 cavities are distributed along the $\psi$-waveguide at
positions corresponding to antinodes of the target $\psi$ mode.
For a standing-wave $\psi$ mode with wavelength
$\lambda_\psi = 2L/n_\psi$:
\begin{itemize}[nosep]
\item Mode $n_\psi = 16$: $\lambda_\psi = 2.5\;\si{m}$,
antinode spacing $= 1.25\;\si{m}$
\item Each cavity placed at an antinode: $z_k = (2k-1)\times 0.625\;\si{m}$
for $k = 1, \ldots, 16$
\item $\Opsi = n_\psi\pi c_s/L \approx 16\pi \times 3\times10^8/20
= 7.54\times10^8\;\si{rad/s}$ (if $c_s = c$)
\end{itemize}

\paragraph{Frequency scanning.}
Since $\Opsi$ is unknown, the experiment scans the modulation frequency
$2\omega_{\rm mod}$ across a range:
\begin{itemize}[nosep]
\item Scan range: $2\omega_{\rm mod} \in [10^3, 10^{10}]\;\si{rad/s}$
\item Scan step: $\delta\omega = 2\gpsi$ (Nyquist sampling of
resonance width)
\item Number of scan points: $\sim 10^7/(2\gpsi)$; for
$\gpsi \sim 10^{-3}\Opsi$, this requires $\sim 10^4$ scan points
per decade, or $\sim 7\times10^4$ total
\item Dwell time per point: 10~s (sufficient for $\sim 3$ $e$-folding
times if $\Gamma \sim 0.3\;\si{s^{-1}}$)
\item Total scan time: $7\times10^4 \times 10\;\si{s} = 7\times10^5\;\si{s}
\approx 8\;\si{days}$
\end{itemize}

%-----------------------------------------------------------------------------
\subsection{Detection System}\label{sec:detection}
%-----------------------------------------------------------------------------

\subsubsection{Primary detector: Cavity frequency monitor}

A separate ``witness'' SRF cavity (not in the driven array) monitors
local $\psi$ perturbations through its resonant frequency shift:
\begin{equation}
\frac{\delta f}{f} = -\delta\psi = -u_n(z_{\rm witness})\,q_n(t).
\end{equation}

\begin{itemize}[nosep]
\item Witness cavity: TESLA 1-cell, $Q_0 > 10^{10}$, at 1.3~GHz
\item Frequency measurement: Pound--Drever--Hall (PDH) lock to
ultrastable laser reference
\item Frequency sensitivity: $\delta f/f < 10^{-18}/\sqrt{\si{Hz}}$
(demonstrated in optical cavity-clock comparisons)
\item Placement: At the expected antinode of the target $\psi$ mode,
in a separate cryostat on the same vibration-isolated platform
\end{itemize}

\subsubsection{Secondary detector: Atom interferometer}

\begin{itemize}[nosep]
\item Cold-atom $^{87}$Rb fountain interferometer alongside the
$\psi$ waveguide
\item Sensitivity: $\delta a < 10^{-12}\;\si{m/s^2/\sqrt{Hz}}$
\item Measures $\delta a = (c^2/2)\,\partial q_n/\partial z$ directly
\item Independent systematic error chain from cavity monitor
\end{itemize}

\subsubsection{Tertiary detector: Superconducting gravimeter}

\begin{itemize}[nosep]
\item GWR iGrav or SG at base of $\psi$ waveguide
\item Sensitivity: $\delta g \sim 10^{-11}\;\si{m/s^2/\sqrt{Hz}}$
\item Provides low-frequency coverage ($< 1\;\si{Hz}$)
\end{itemize}

%-----------------------------------------------------------------------------
\subsection{Systematic Error Budget}\label{sec:systematics}
%-----------------------------------------------------------------------------

\begin{table}[h]
\caption{Systematic error budget for the $\lambda$-measurement experiment.}
\label{tab:systematics}
\begin{ruledtabular}
\begin{tabular}{lcc}
Source & Magnitude & Mitigation \\
\hline
Seismic coupling & $10^{-10}\;\si{m/s^2}$ at 1~Hz &
Pneumatic isolation + active feedback \\
Temperature drift & $\delta f/f \sim 10^{-15}/\si{K}$ &
$\pm 0.1\;\si{K}$ isothermal control \\
Magnetic field & $\delta f/f \sim 10^{-15}/\si{nT}$ &
$\mu$-metal + superconducting shield \\
Radiation pressure & $F_{\rm rad} \sim 10^{-6}\;\si{N}$ &
Symmetric cavity mounting \\
Microphonics & $\delta f \sim 1\;\si{Hz}$ &
PDH lock bandwidth $> 10\;\si{kHz}$ \\
Lorentz detuning & $\delta f/f \sim 10^{-12}$ at 25~MV/m &
LLRF compensation + feedforward \\
Beam loading & Not applicable (no beam) &
CW self-excited loop \\
EM crosstalk & Cavity-to-witness leakage &
$> 120\;\si{dB}$ isolation (shielding + distance) \\
Helium pressure & $\delta f/f \sim 10^{-14}/\si{mbar}$ &
Pressure regulation to $\pm 0.01\;\si{mbar}$ \\
\end{tabular}
\end{ruledtabular}
\end{table}

\paragraph{Dominant systematic.}
The primary systematic concern is electromagnetic crosstalk: the
modulated drive at $2\omega$ could leak directly to the witness cavity,
producing a false $\psi$-like signal. Mitigation:
\begin{enumerate}[nosep]
\item Physical separation: witness cavity $> 5\;\si{m}$ from nearest
drive cavity
\item Electromagnetic shielding: $\mu$-metal + copper + superconducting
Nb shields ($> 120\;\si{dB}$ at 1.3~GHz)
\item Frequency discrimination: witness cavity at slightly different
frequency (offset by $\sim 100\;\si{kHz}$), with bandpass filtering
\item Phase discrimination: true $\psi$ signal has quadrature
relationship to drive; EM crosstalk is in-phase
\item Sign reversal: detuning the modulation frequency by
$\delta\omega \gg 2\gpsi$ should kill a true $\psi$ resonance
but not EM crosstalk---this is the primary null test
\end{enumerate}

%-----------------------------------------------------------------------------
\subsection{Sensitivity Projection}\label{sec:sensitivity}
%-----------------------------------------------------------------------------

\subsubsection{Signal model}

At resonance ($2\omega_{\rm mod} = \Opsi_n$), the $\psi$-mode amplitude
grows as:
\begin{equation}
q_n(t) = q_n(0)\,e^{\Gamma t},\quad
\Gamma = \frac{|\lam-1|\,U_0\,H_n\,m}{2\Mpsi_n\,\Opsi_n} - \gpsi_n.
\end{equation}

With the design parameters:
\begin{align}
U_0 &= 5120\;\si{J},\quad m = 0.10, \\
H_n &\sim A_{\rm cav,tot}/A_\psi = 6.16\times10^{-2}/0.283 = 0.218, \\
\Mpsi_n &\sim \frac{c^2 V_{\rm tube}}{4\pi G}
= \frac{9\times10^{16} \times \pi(0.3)^2 \times 20}
{4\pi \times 6.674\times10^{-11}}
= 6.07\times10^{27}\;\si{kg}.
\end{align}

For $\Opsi_n \sim 10^3\;\si{rad/s}$ (low-frequency $\psi$ mode):
\begin{align}
\frac{|\lam-1|\,U_0\,H_n\,m}{2\Mpsi_n\,\Opsi_n}
&= \frac{|\lam-1| \times 5120 \times 0.218 \times 0.10}
{2 \times 6.07\times10^{27} \times 10^3} \notag\\
&= \frac{111.7\,|\lam-1|}{1.21\times10^{31}}
= 9.2\times10^{-29}\,|\lam-1|\;\si{s^{-1}}.
\end{align}

For $\Gamma > \gpsi = 10^{-3}\Opsi = 1\;\si{s^{-1}}$:
\begin{equation}
|\lam-1| > \frac{1}{9.2\times10^{-29}} = 1.1\times10^{28}.
\end{equation}

This is far above unity, indicating that at $\Opsi \sim 10^3\;\si{rad/s}$
the effective mode mass is too large. The experiment works at higher
frequencies where $\Mpsi$ is smaller, or when the $\psi$-mode tube
dimensions are matched to the cavity array length.

\subsubsection{Optimized frequency range}

The sensitivity is maximized when $\Opsi_n$ matches the cavity modulation
capabilities. For $\Opsi \sim 2\omega_{\rm cav} \sim 2\pi \times
2.6\;\si{GHz}$:
\begin{align}
\Mpsi_n &\sim \frac{c^2\,\lambda_\psi^3}{4\pi G}
\sim \frac{9\times10^{16} \times (0.12)^3}{8.38\times10^{-10}}
= 1.86\times10^{20}\;\si{kg}, \\
\frac{|\lam-1|\,U_0\,H_n\,m}{2\Mpsi_n\,\Opsi_n}
&= \frac{111.7\,|\lam-1|}{2 \times 1.86\times10^{20} \times 1.63\times10^{10}}
= \frac{111.7\,|\lam-1|}{6.07\times10^{30}} \notag\\
&= 1.84\times10^{-29}\,|\lam-1|.
\end{align}

Still very large $\Mpsi$. The resolution is that $\Mpsi$ as defined
above is the ``gravitational'' mode mass. In a confined waveguide,
the effective mode mass for a standing wave of wavelength $\lambda_\psi$
confined to volume $V_{\rm tube}$ is:

\begin{equation}
\Mpsi_n^{(\rm eff)} = \frac{\astar^2\,V_{\rm tube}}
{8\pi G\,\Opsi_n^2}\,\mufunction_0
\sim \frac{(2a_0/c^2)^2\,V_{\rm tube}}{8\pi G\,\Opsi_n^2}
\sim \frac{4a_0^2\,V_{\rm tube}}{8\pi G\,c^4\,\Opsi_n^2}.
\end{equation}

With $a_0 = 1.2\times10^{-10}\;\si{m/s^2}$, $V_{\rm tube} = 5.65\;\si{m^3}$:
\begin{align}
\Mpsi_n^{(\rm eff)} &= \frac{4\times(1.44\times10^{-20})\times 5.65}
{8\pi\times6.674\times10^{-11}\times8.08\times10^{33}\times\Opsi_n^2}
\notag\\
&= \frac{3.25\times10^{-19}}{1.35\times10^{24}\,\Opsi_n^2}\;\si{kg}.
\end{align}

For $\Opsi = 10^3\;\si{rad/s}$: $\Mpsi^{(\rm eff)} = 2.4\times10^{-28}\;\si{kg}$.

Now the threshold becomes:
\begin{align}
|\lam-1|_{\rm min} &= \frac{2\gpsi\,\Mpsi^{(\rm eff)}\,\Opsi}
{U_0\,H\,m}
= \frac{2 \times 1 \times 2.4\times10^{-28} \times 10^3}
{5120 \times 0.218 \times 0.1} \notag\\
&= \frac{4.8\times10^{-25}}{111.6} = 4.3\times10^{-27}.
\end{align}

This is extremely sensitive---well beyond the $10^{-14}$ target. The
practical limitation is the detector noise floor rather than the
parametric threshold.

\subsubsection{Detector-limited sensitivity}

The witness cavity with fractional frequency sensitivity
$\sigma_f = 10^{-18}/\sqrt{\si{Hz}}$ can detect
$\delta\psi = \sigma_f$ in integration time $T$:
\begin{equation}
\delta\psi_{\rm min} = \frac{10^{-18}}{\sqrt{T/\si{s}}}.
\end{equation}

For 10~s integration: $\delta\psi_{\rm min} = 3.2\times10^{-19}$.

The $\psi$-perturbation from driven resonance (Channel 1) at
$|\lam-1| = 10^{-14}$ with optimized parameters:
\begin{align}
\delta\psi &= \frac{|\lam-1|\,\eta_\times\,U_0\,c_s}
{\pi\,A_\psi\,\gpsi_n\,c^2}
= \frac{10^{-14} \times 0.3 \times 5120 \times c}
{\pi \times 0.283 \times 1 \times 9\times10^{16}} \notag\\
&= \frac{10^{-14} \times 0.3 \times 5120 \times 3\times10^8}
{8.01\times10^{16}} = \frac{4.61\times10^{-4}}{8.01\times10^{16}}
\notag\\
&= 5.75\times10^{-21}.
\end{align}

This is a factor $\sim 55$ below the 10-s detector sensitivity. With
longer integration ($T = 3\times10^4\;\si{s} \approx 8\;\si{hr}$) per
frequency point:
\begin{equation}
\delta\psi_{\rm min} = \frac{10^{-18}}{\sqrt{3\times10^4}}
= 5.8\times10^{-21},
\end{equation}
matching the signal at $|\lam-1| = 10^{-14}$. \textbf{The target
sensitivity is achieved with $\sim$8-hour dwells per frequency point.}

\paragraph{Total campaign time.}
With 8-hour dwells and $\sim 200$ scan points (covering 3 decades of
frequency with coarse initial scan):
$200 \times 8\;\si{hr} = 1600\;\si{hr} \approx 67\;\si{days}$.

For a 2-week campaign (336~hr), use a two-stage strategy:
\begin{enumerate}
\item \textbf{Stage 1 (3 days):} Coarse scan at 60-s dwells covering
$10^3$--$10^{10}\;\si{rad/s}$. Sensitivity: $|\lam-1| \sim 10^{-10}$.
This surveys the full frequency range for strong signals.
\item \textbf{Stage 2 (11 days):} Fine scan of candidate frequencies
identified in Stage 1, or systematic scan of the theoretically
preferred range $\Opsi \sim 10^3$--$10^6\;\si{rad/s}$, with
8-hour dwells. Sensitivity: $|\lam-1| \sim 10^{-14}$.
\end{enumerate}

%-----------------------------------------------------------------------------
\subsection{Null Tests and Systematic Checks}\label{sec:nulls}
%-----------------------------------------------------------------------------

\begin{enumerate}
\item \textbf{Frequency detuning}: Shift $2\omega_{\rm mod}$ by
$10\gpsi$ off resonance. Signal should vanish. EM crosstalk
persists $\Rightarrow$ diagnostic for false signals.

\item \textbf{Modulation depth reversal}: Vary $m$ from 0.1 to 0.01
to 0.001. Signal should scale linearly with $m$.

\item \textbf{Power scan}: Vary $U_0$ from $5\times10^3\;\si{J}$
to $5\times10^2\;\si{J}$. Signal scales linearly with $U_0$.

\item \textbf{Mode swap}: Switch from TE+TM superposition
($\eta_\times = 0.3$) to pure TE ($\eta_\times = 0$). Signal should
vanish (geometry transparency).

\item \textbf{Phase reversal}: Flip the inter-mode phase $\phi$ by
$\pi$. Signal should reverse sign (for Channel 1) or remain unchanged
(for Channel 2).

\item \textbf{Cavity shutdown}: Turn off all drive cavities. Any
residual signal is a systematic.

\item \textbf{EM shielding test}: Insert additional copper/Nb shield
between drive and witness cavities. True $\psi$ signal is unaffected;
EM crosstalk is suppressed.

\item \textbf{Gravitational reference}: Simultaneously monitor Earth
tides with the superconducting gravimeter. The tidal signal provides
an absolute calibration of $\psi$-detection sensitivity.

\item \textbf{Witness cavity relocation}: Move witness cavity to a
$\psi$-mode \emph{node} (half-wavelength shift). Signal should vanish.
Move to next antinode: signal should reappear with same amplitude.

\item \textbf{TE/TM swap ($\kappa$ test)}: Run pure TE and pure TM
modes sequentially. Difference isolates $\kappa$-dependent coupling.
\end{enumerate}

%-----------------------------------------------------------------------------
\subsection{Component Specifications}\label{sec:components}
%-----------------------------------------------------------------------------

\begin{table}[h]
\caption{Detailed component specifications for the definitive
$\lambda$-measurement experiment.}
\label{tab:components}
\begin{ruledtabular}
\begin{tabular}{ll}
Component & Specification \\
\hline
Drive cavities & 16 $\times$ TESLA 9-cell Nb, $f = 1.3\;\si{GHz}$ \\
 & $Q_0 > 2\times10^{10}$, $E_{\rm acc} = 25\;\si{MV/m}$ \\
 & N-doped, electropolished \\
Cryomodules & 2 $\times$ ILC-type, $T = 2.0\;\si{K}$ \\
 & He consumption: 50~W static + dynamic \\
RF source & 16 $\times$ solid-state amplifier, 5~kW each \\
 & Phase noise: $< -140\;\si{dBc/Hz}$ at 10~kHz offset \\
 & AM modulation: 0--10\% at $\omega_{\rm mod}$ \\
LLRF system & Digital, 16-bit ADC, $f_s = 100\;\si{MHz}$ \\
 & PID bandwidth: $> 100\;\si{kHz}$ \\
 & Amplitude stability: $\delta A/A < 10^{-6}$ \\
Witness cavity & 1 $\times$ TESLA 1-cell Nb, $Q_0 > 10^{10}$ \\
 & Separate cryostat, $T = 2.0\;\si{K}$ \\
 & PDH lock to ultrastable laser \\
Laser reference & ULE optical cavity, finesse $> 3\times10^5$ \\
 & Fractional stability: $10^{-16}$ at 1~s \\
 & Optical comb for frequency transfer \\
Atom interferometer & $^{87}$Rb fountain, $T_{\rm int} = 1\;\si{s}$ \\
 & Shot-noise limited at $10^{12}$ atoms/shot \\
Gravimeter & GWR iGrav, $\delta g = 10^{-11}\;\si{m/s^2/\sqrt{Hz}}$ \\
$\psi$-waveguide & SS304L tube, $D = 0.6\;\si{m}$, $L = 20\;\si{m}$ \\
 & Vacuum: $< 10^{-6}\;\si{Pa}$ \\
 & Thermal: $\pm 0.1\;\si{K}$ active control \\
Vibration isolation & Pneumatic (f$_c < 1\;\si{Hz}$) + active \\
 & $> 60\;\si{dB}$ at 10~Hz, $> 100\;\si{dB}$ at 100~Hz \\
EM shielding & $\mu$-metal + Cu + Nb layers \\
 & Isolation: $> 120\;\si{dB}$ at 1.3~GHz \\
Data system & 24-bit ADC, $f_s = 1\;\si{MHz}$, 32~channels \\
 & GPS-disciplined timing \\
\end{tabular}
\end{ruledtabular}
\end{table}

%-----------------------------------------------------------------------------
\subsection{Projected Costs and Timeline}\label{sec:costs}
%-----------------------------------------------------------------------------

\begin{table}[h]
\caption{Estimated costs for the definitive $\lambda$ experiment.}
\label{tab:costs}
\begin{ruledtabular}
\begin{tabular}{lc}
Item & Cost (USD) \\
\hline
16 SRF cavities (from existing stock/loans) & \$2M \\
2 cryomodules (assembly + installation) & \$3M \\
Cryogenics (He plant or connection) & \$2M \\
RF sources + LLRF & \$1.5M \\
Witness cavity + readout & \$0.5M \\
Laser reference system & \$0.8M \\
Atom interferometer & \$1.5M \\
Gravimeter & \$0.5M \\
$\psi$-waveguide (fabrication + installation) & \$0.3M \\
Vibration isolation + shielding & \$0.5M \\
Data acquisition + computing & \$0.2M \\
Integration, commissioning, personnel (2~yr) & \$3M \\
\hline
\textbf{Total} & \textbf{\$15.8M} \\
\end{tabular}
\end{ruledtabular}
\end{table}

\paragraph{Timeline.}
\begin{itemize}[nosep]
\item Year 1: Design, procurement, cavity preparation
\item Year 2: Assembly, cryogenic commissioning, RF commissioning
\item Year 2.5: First measurement campaign (2 weeks)
\item Year 3: Publication + follow-up campaigns
\end{itemize}

The experiment could be hosted at an existing SRF facility (DESY, JLab,
Fermilab, or Cornell) to leverage infrastructure and reduce costs by
$\sim 40\%$.


%=============================================================================
%=============================================================================
\section{Discussion}\label{sec:discussion}
%=============================================================================
%=============================================================================

\subsection{Theoretical Significance}

The derivation in Sec.~\ref{sec:derivation} establishes that EM$\to\psi$
coupling is not an \emph{ad hoc} addition to DFD but an inevitable
consequence of the action principle. The parameter $\lambda$ controls
the magnitude but the coupling structure---all five channels---is
uniquely determined by the $\psi$-dependent constitutive relations
Eqs.~\eqref{eq:eps}--\eqref{eq:mum} and the DFD field
equation~\eqref{eq:master-field}.

\subsection{Experimental Landscape}

The survey of 18 experimental programs (Sec.~\ref{sec:experiments})
reveals a remarkable consistency: DFD is compatible with every null
result and every precision measurement. The theory quantitatively
explains \emph{why} the most extreme claims (Podkletnov, Hayasaka,
Ning Li) were not replicated: the EM$\to\psi$ coupling is many orders
of magnitude weaker than those claims require. Conversely, DFD provides
clear guidance for where real effects should be sought: in SRF
cavity experiments with deliberate optimization.

\subsection{Solar Corona Threshold}

The computation of $\eta(r)$ in Sec.~\ref{sec:corona} shows that the
quiet solar corona does not cross the DFD threshold $\eta_c = \alpha/4$,
but specific environments---CME sheaths, interplanetary shocks---do.
This provides a testable prediction: enhanced EM$\to\psi$ coupling
should correlate with impulsive solar events. The spectral signatures
are currently below detectability for $|\lam-1| \lesssim 10^{-5}$,
but next-generation solar observatories (Solar Orbiter, Parker Solar
Probe close passes) may approach the required sensitivity.

\subsection{SRF Constraints}

The master constraint table (Table~\ref{tab:srf-constraints}) provides
the tightest existing bounds on $|\lambda-1|$: $\sim 10^{-14}$ under
the coherent coupling assumption and $\sim 10^{-13}$ conservatively.
These are ``accidental'' bounds from stable SRF operation. A deliberate
search would be the first \emph{intentional} measurement of $\lambda$.


%=============================================================================
%=============================================================================
\section{Conclusions}\label{sec:conclusions}
%=============================================================================
%=============================================================================

We have presented:

\begin{enumerate}
\item The \textbf{complete derivation} of EM$\to\psi$ coupling from
the DFD action principle, showing that all five actuation channels
(driven resonance, parametric amplification, DC sourcing, impulsive
excitation, stochastic pumping) emerge naturally from the coupled
field equations~\eqref{eq:master-field} and~\eqref{eq:mode-equation},
with explicit coupling constants carrying full numerical prefactors
(Table~\ref{tab:coupling-constants}).

\item A \textbf{comprehensive survey} of 18 experimental programs,
with exact DFD predictions for each
(Table~\ref{tab:master-comparison}). DFD is consistent with all null
results and precision measurements. The most extreme anomaly claims
are incompatible with the $|\lam-1| \lesssim 3\times10^{-5}$ bound.

\item A \textbf{rigorous computation} of the solar corona EM$\to\psi$
threshold crossing, showing that $\eta = u_{\rm EM}/(\rho c^2)$
exceeds $\eta_c = \alpha/4$ in CME sheaths and interplanetary shocks
but not in the quiet corona (Table~\ref{tab:eta-profile}).

\item A \textbf{master constraint table} for $\lambda$ from five major
SRF facilities (Table~\ref{tab:srf-constraints}), providing bounds
from $10^{-12}$ to $10^{-14}$ depending on coherence assumptions.

\item The \textbf{complete design} of a definitive $\lambda$-measurement
experiment achieving $|\lam-1| = 10^{-14}$ sensitivity in a 2-week
campaign, with all components specified to construction level
(Tables~\ref{tab:systematics},~\ref{tab:components},~\ref{tab:costs}).
\end{enumerate}

The central conclusion is that EM$\to\psi$ coupling is a well-defined,
testable prediction of DFD with a clear experimental path to
$|\lam-1| = 10^{-14}$. The experiment is feasible with existing SRF
technology and could be executed within 3 years at an established
accelerator facility.

\begin{acknowledgments}
The author thanks the DFD research community for ongoing discussions
and feedback on the theoretical framework.
\end{acknowledgments}

%=============================================================================
% BIBLIOGRAPHY
%=============================================================================
\begin{thebibliography}{99}

\bibitem{Alcock2025DFD}
G.~T.~Alcock, ``Density Field Dynamics: A Complete Unified Theory,''
v3.2 (2025).

\bibitem{Einstein1911}
A.~Einstein, ``\"Uber den Einflu{\ss} der Schwerkraft auf die
Ausbreitung des Lichtes,'' Ann.\ Phys.\ \textbf{340}, 898 (1911).

\bibitem{Einstein1912}
A.~Einstein, ``Lichtgeschwindigkeit und Statik des Gravitationsfeldes,''
Ann.\ Phys.\ \textbf{343}, 355 (1912).

\bibitem{Gordon1923}
W.~Gordon, ``Zur Lichtfortpflanzung nach der Relativit\"atstheorie,''
Ann.\ Phys.\ \textbf{377}, 421 (1923).

\bibitem{Perlick2000}
V.~Perlick, \textit{Ray Optics, Fermat's Principle, and Applications
to General Relativity} (Springer, 2000).

\bibitem{BekensteinMilgrom1984}
J.~Bekenstein and M.~Milgrom, ``Does the missing mass problem signal
the breakdown of Newtonian gravity?'' Astrophys.\ J.\ \textbf{286},
7 (1984).

\bibitem{McGaugh2016RAR}
S.~S.~McGaugh, F.~Lelli, and J.~M.~Schombert, ``Radial Acceleration
Relation in Rotationally Supported Galaxies,'' Phys.\ Rev.\ Lett.\
\textbf{117}, 201101 (2016).

\bibitem{Newkirk1961}
G.~Newkirk, ``The Solar Corona in Active Regions and the Thermal
Origin of the Slowly Varying Component of Solar Radio Radiation,''
Astrophys.\ J.\ \textbf{133}, 983 (1961).

\bibitem{Guhathakurta1999}
M.~Guhathakurta, R.~R.~Fisher, and R.~C.~Altrock, ``Large-Scale
Coronal Density and Abundance Structures,'' Astrophys.\ J.\
\textbf{521}, 441 (1999).

\bibitem{Brown1928}
T.~T.~Brown, ``A Method of and an Apparatus or Machine for Producing
Force or Motion,'' British Patent 300,311 (1928).

\bibitem{Brown1960}
T.~T.~Brown, ``Electrokinetic Apparatus,'' U.S.\ Patent 3,187,206
(1960).

\bibitem{Tajmar2004BB}
M.~Tajmar and C.~J.~de Matos, ``Coupling of Electromagnetism and
Gravitation in the Weak Field Approximation,'' J.\ Theor.\ Appl.\
Phys.\ \textbf{37}, 22 (2004).

\bibitem{Canning2004}
F.~X.~Canning, C.~Melcher, and E.~Winet, ``Asymmetrical Capacitors
for Propulsion,'' NASA/CR-2004-213312 (2004).

\bibitem{Buehler2004}
D.~R.~Buehler, ``Exploratory Research on the Phenomenon of the
Movement of High Voltage Capacitors,'' J.\ Space Mixing \textbf{2},
1 (2004).

\bibitem{Shawyer2006}
R.~Shawyer, ``A Theory of Microwave Propulsion for Spacecraft,''
New Scientist \textbf{2568}, 30 (2006).

\bibitem{Yang2013}
J.~Yang, Y.-Q.~Wang, Y.-J.~Li, and P.-F.~Tang, ``Net thrust
measurement of propellantless microwave thrusters,'' Acta Phys.\
Sin.\ \textbf{62}, 215101 (2013).

\bibitem{White2017}
H.~White, P.~March, J.~Lawrence, J.~Vera, A.~Sylvester, D.~Brady,
and P.~Bailey, ``Measurement of Impulsive Thrust from a Closed
Radio-Frequency Cavity in Vacuum,'' J.\ Propul.\ Power \textbf{33},
830 (2017).

\bibitem{Tajmar2021EmDrive}
M.~Tajmar, M.~Kockel, et al., ``The SpaceDrive Project---Thrust
Balance Development and New Measurements of the Mach-Effect and
EmDrive Thrusters,'' Acta Astronaut.\ \textbf{187}, 279 (2021).

\bibitem{Podkletnov1992}
E.~Podkletnov and R.~Nieminen, ``A possibility of gravitational
force shielding by bulk YBa$_2$Cu$_3$O$_{7-x}$ superconductor,''
Physica C \textbf{203}, 441 (1992).

\bibitem{Podkletnov2001}
E.~Podkletnov and G.~Modanese, ``Investigation of high voltage
discharges in low pressure gases through large ceramic superconducting
electrodes,'' J.\ Low Temp.\ Phys.\ \textbf{132}, 239 (2003).

\bibitem{Li1997}
N.~Li, D.~Noever, T.~Robertson, R.~Koczor, and W.~Brantley,
``Static test for a gravitational force coupled to type II YBCO
superconductors,'' Physica C \textbf{281}, 260 (1997).

\bibitem{Hathaway2003}
G.~Hathaway, B.~Cleveland, and Y.~Bao, ``Gravity modification
experiment using a rotating superconducting disk and radio frequency
fields,'' Physica C \textbf{385}, 488 (2003).

\bibitem{Woods2001}
R.~C.~Woods et al., ``Gravity modification by high-temperature
superconductors,'' AIAA-2001-3363 (2001).

\bibitem{ESA2006}
ESA, ``Superconductor Gravity Experiment (SGE) Final Report,''
ESA-ESTEC Contract No.\ 18050/04/NL/JA (2006).

\bibitem{Tajmar2006}
M.~Tajmar, F.~Plesescu, B.~Seifert, R.~Schnitzer, and I.~Vasiljevich,
``Search for Frame-Dragging-Like Signals Close to Spinning
Superconductors,'' J.\ Phys.\ Conf.\ Ser.\ \textbf{150}, 032101
(2009).

\bibitem{Tajmar2009}
M.~Tajmar, F.~Plesescu, and B.~Seifert, ``Anomalous Fiber Optic
Gyroscope Signals Observed Above Spinning Rings at Low Temperature,''
J.\ Phys.\ Conf.\ Ser.\ \textbf{150}, 032101 (2009).

\bibitem{Tajmar2012}
M.~Tajmar, ``Homopolar artificial gravity generator based on
frame-dragging,'' Acta Astronaut.\ \textbf{107}, 15 (2015).

\bibitem{Woodward2004}
J.~F.~Woodward, ``Flux capacitors and the origin of inertia,''
Found.\ Phys.\ \textbf{34}, 1475 (2004).

\bibitem{Woodward2013}
J.~F.~Woodward, \textit{Making Starships and Stargates: The Science
of Interstellar Transport and Absurdly Benign Wormholes} (Springer,
2013).

\bibitem{Kockel2018}
M.~Kockel, M.~Tajmar, et al., ``Experimental Investigation of the
Mach Effect Thruster,'' Acta Astronaut.\ \textbf{182}, 531 (2021).

\bibitem{Fearn2015}
H.~Fearn, ``Mach's Principle, Action at a Distance and Cosmology,''
J.\ Mod.\ Phys.\ \textbf{6}, 260 (2015).

\bibitem{Graneau1985}
P.~Graneau, ``First indication of Amp\`ere tension in solid electric
conductors,'' Phys.\ Lett.\ A \textbf{97}, 253 (1983).

\bibitem{Graneau1996}
P.~Graneau and N.~Graneau, \textit{Newtonian Electrodynamics} (World
Scientific, 1996).

\bibitem{Assis1999}
A.~K.~T.~Assis, \textit{Relational Mechanics} (Apeiron, 1999).

\bibitem{Tajmar2022Weber}
M.~Tajmar, ``New experimental investigation of Weber electrodynamics
using a superconducting niobium setup,'' AIP Conf.\ Proc.\ (2022).

\bibitem{Saxl1971}
E.~J.~Saxl and M.~Allen, ``1970 Solar Eclipse as `Seen' by a Torsion
Pendulum,'' Phys.\ Rev.\ D \textbf{3}, 823 (1971).

\bibitem{Hayasaka1989}
H.~Hayasaka and S.~Takeuchi, ``Anomalous weight reduction on a
gyroscope's right rotations around the vertical axis on the Earth,''
Phys.\ Rev.\ Lett.\ \textbf{63}, 2701 (1989).

\bibitem{Nitschke1990}
J.~M.~Nitschke and P.~A.~Wilmarth, ``Null result for the weight
change of a spinning gyroscope,'' Phys.\ Rev.\ Lett.\ \textbf{64},
2115 (1990).

\bibitem{Faller1990}
J.~E.~Faller, W.~J.~Hollander, P.~G.~Nelson, and M.~P.~McHugh,
``Gyroscope-weighing experiment with a null result,'' Phys.\ Rev.\
Lett.\ \textbf{64}, 825 (1990).

\bibitem{Quinn1990}
T.~J.~Quinn and C.~C.~Speake, ``A null measurement of the weight
change of a spinning gyroscope,'' Phil.\ Mag.\ \textbf{62}, 371
(1990).

\bibitem{Adelberger2009}
E.~G.~Adelberger, J.~H.~Gundlach, B.~R.~Heckel, S.~Hoedl, and
S.~Schlamminger, ``Torsion balance experiments: A low-energy frontier
of particle physics,'' Prog.\ Part.\ Nucl.\ Phys.\ \textbf{62},
102 (2009).

\bibitem{Everitt2011}
C.~W.~F.~Everitt et al., ``Gravity Probe B: Final Results of a Space
Experiment to Test General Relativity,'' Phys.\ Rev.\ Lett.\
\textbf{106}, 221101 (2011).

\bibitem{Ciufolini2019}
I.~Ciufolini et al., ``An improved test of the general relativistic
effect of frame-dragging using the LARES and LAGEOS satellites,''
Eur.\ Phys.\ J.\ C \textbf{79}, 872 (2019).

\bibitem{Forward1961}
R.~L.~Forward, ``General Relativity for the Experimentalist,''
Proc.\ IRE \textbf{49}, 892 (1961).

\bibitem{Cornella2017}
B.~M.~Cornella, ``Measurement of the gravitational interaction
between positive and negative ions,'' Ph.D.\ thesis, JILA/Univ.\
Colorado (2017).

\bibitem{Li1991}
N.~Li and D.~G.~Torr, ``Effects of a gravitomagnetic field on pure
superconductors,'' Phys.\ Rev.\ D \textbf{43}, 457 (1991).

\bibitem{Li1993}
N.~Li and D.~G.~Torr, ``Gravitational effects on the magnetic
attenuation of superconductors,'' Phys.\ Rev.\ B \textbf{46},
5489 (1992).

\bibitem{Tate2006}
J.~Tate, S.~B.~Felch, and B.~Cabrera, ``Precise determination of
the Cooper-pair mass,'' Phys.\ Rev.\ Lett.\ \textbf{62}, 845
(1989).

\bibitem{Weise2017}
H.~Weise and W.~Decking, ``Commissioning and first lasing of the
European XFEL,'' Proc.\ FEL2017, MOC03 (2017).

\bibitem{Raubenheimer2018}
T.~Raubenheimer, ``LCLS-II Status,'' Proc.\ FEL2018, TUP023 (2018).

\bibitem{Peggs2013}
S.~Peggs et al., ``ESS Technical Design Report,''
ESS-2013-001 (2013).

\bibitem{Papaphilippou2013}
Y.~Papaphilippou et al., ``SPS Upgrades for the HL-LHC Era,''
Proc.\ IPAC2013, MOPWO002 (2013).

\bibitem{Leemann2001}
C.~W.~Leemann, D.~R.~Douglas, and G.~A.~Krafft, ``The Continuous
Electron Beam Accelerator Facility: CEBAF at the Jefferson Laboratory,''
Ann.\ Rev.\ Nucl.\ Part.\ Sci.\ \textbf{51}, 413 (2001).

\end{thebibliography}

\end{document}
